
\documentclass[11pt,a4paper]{article}

\usepackage{iftex}
\ifPDFTeX
  \usepackage[T2A]{fontenc}
  \usepackage[utf8]{inputenc}
  \usepackage[russian]{babel}
\else
  \usepackage{fontspec}
  \usepackage[russian]{babel}
  \setmainfont{DejaVu Serif}
  \setsansfont{DejaVu Sans}
  \setmonofont{DejaVu Sans Mono}
\fi
\usepackage{amsmath,amssymb,mathtools,bm}
\usepackage{geometry}
\usepackage{enumitem}
\usepackage{hyperref}
\usepackage{booktabs}
\usepackage{array}
\geometry{margin=1.8cm}
\setlist{nosep,leftmargin=*}
\hypersetup{colorlinks=true,linkcolor=blue,urlcolor=blue}

\newcommand{\E}{\mathbb{E}}
\renewcommand{\Prob}{\mathbb{P}}
\newcommand{\R}{\mathbb{R}}
\newcommand{\KL}{\mathrm{KL}}
\newcommand{\Var}{\mathrm{Var}}
\newcommand{\argmax}{\operatorname*{arg\,max}}
\newcommand{\argmin}{\operatorname*{arg\,min}}
\newcommand{\softmax}{\operatorname{softmax}}
\newcommand{\clip}{\operatorname{clip}}
\newcommand{\sg}{\operatorname{sg}}
\newcommand{\1}{\mathbf{1}}
\newcommand{\cS}{\mathcal S}
\newcommand{\cA}{\mathcal A}
\newcommand{\cD}{\mathcal D}
\newcommand{\cN}{\mathcal N}
\newcommand{\cT}{\mathcal T}
\newcommand{\dd}{\,d}

\title{Краткие TeX-заготовки ответов к экзамену по RL}
\author{по программе экзамена и лекциям курса}
\date{}

\begin{document}
\maketitle
\tableofcontents
\newpage

\section{Кросс-энтропийный метод: общий вид, оптимизация и RL}

\subsection*{Идея}
Кросс-энтропийный метод (Cross-Entropy Method, CEM) можно понимать как метод поиска хорошего распределения над решениями. Мы не оптимизируем одну точку напрямую, а задаем параметрическое семейство распределений $q_\theta(x)$ над кандидатами $x$, сэмплируем кандидатов, отбираем элиту и подгоняем распределение так, чтобы оно с большей вероятностью порождало элитные решения.

Общая задача оптимизации:
\[
    x^* \in \argmax_{x\in \mathcal X} S(x),
\]
где $S(x)$ - целевая функция. Вместо прямого поиска $x^*$ вводим распределение $q_\theta(x)$ и итеративно сдвигаем его массу в сторону областей, где $S(x)$ велико.

\subsection*{Редкое событие и условное распределение на успехе}
Исторически CEM связан с оценкой вероятности редких событий. Пусть $X\sim p(x)$ и нас интересует событие
\[
    S(X) \ge \gamma.
\]
Идеальное распределение для importance sampling - это исходное распределение, условленное на успех:
\[
    p^*(x) = p(x\mid S(x)\ge \gamma)
    = \frac{p(x)\1\{S(x)\ge \gamma\}}{\Prob_p(S(X)\ge \gamma)}.
\]
Здесь $\1\{\cdot\}$ - индикатор: он равен $1$, если условие выполнено, и $0$ иначе. Такое распределение идеально, потому что оно сэмплирует только успешные объекты. На практике $p^*$ неизвестно, поэтому его приближают распределением $q_\theta$ из удобного семейства.

Параметры выбирают через минимизацию KL-дивергенции:
\[
    \theta^* = \argmin_\theta \KL(p^*\|q_\theta).
\]
Раскрывая KL,
\[
    \KL(p^*\|q_\theta)
    = \E_{p^*}\left[\log \frac{p^*(X)}{q_\theta(X)}\right]
    = \E_{p^*}[\log p^*(X)] - \E_{p^*}[\log q_\theta(X)].
\]
Первое слагаемое не зависит от $\theta$, поэтому задача эквивалентна максимизации правдоподобия элитных сэмплов:
\[
    \theta^* = \argmax_\theta \E_{p^*}\log q_\theta(X).
\]
Отсюда возникает логарифм: это обычная MLE-задача для распределения, которое должно порождать успешные примеры.

\subsection*{Практический алгоритм CEM}
Поскольку $p^*$ неизвестно, его заменяют эмпирическим распределением на элитных сэмплах.

На итерации $k$:
\begin{enumerate}
    \item Имеем текущее распределение $q_{\theta_k}(x)$.
    \item Сэмплируем кандидатов:
    \[
        x_1,\dots,x_N \sim q_{\theta_k}(x).
    \]
    \item Считаем оценки качества $S(x_i)$.
    \item Выбираем порог элиты $\gamma_k$, например $(1-\rho)$-квантиль по значениям $S(x_i)$, где $\rho\in(0,1)$ - доля элитных объектов.
    \item Формируем элитное множество:
    \[
        \mathcal E_k = \{x_i: S(x_i)\ge \gamma_k\}.
    \]
    \item Обновляем параметры распределения по MLE:
    \[
        \theta_{k+1} = \argmax_\theta \sum_{x_i\in\mathcal E_k}\log q_\theta(x_i).
    \]
\end{enumerate}

Если $q_\theta$ - нормальное распределение $\mathcal N(\mu,\Sigma)$, то обновление имеет простой вид:
\[
    \mu_{k+1}=\frac1{|\mathcal E_k|}\sum_{x_i\in\mathcal E_k}x_i,
    \qquad
    \Sigma_{k+1}=\frac1{|\mathcal E_k|}\sum_{x_i\in\mathcal E_k}(x_i-\mu_{k+1})(x_i-\mu_{k+1})^\top.
\]
Часто используют сглаживание:
\[
    \theta_{k+1}\leftarrow \alpha\theta_{k+1}^{\text{MLE}}+(1-\alpha)\theta_k,
\]
чтобы распределение не схлопнулось слишком быстро.

\subsection*{CEM как метод оптимизации}
В детерминированной оптимизации событие успеха задают через высокий уровень целевой функции. Тогда CEM можно записать как итеративное приближение распределения
\[
    q_{k+1}(x) \approx q_k(x\mid S(x)\ge \gamma_k).
\]
Это интуитивно: берем текущее распределение, смотрим, какие сэмплы оказались лучшими, и обучаем новое распределение быть похожим на распределение лучших сэмплов.

В отличие от градиентных методов, CEM:
\begin{itemize}
    \item не требует дифференцируемости $S(x)$;
    \item естественно работает с дискретными, непрерывными и смешанными переменными;
    \item может быть устойчив к локальной негладкости;
    \item но требует много сэмплов и страдает в больших размерностях.
\end{itemize}

\subsection*{CEM в обучении с подкреплением}
В RL объектом оптимизации может быть не одно действие, а политика или целая траектория. Пусть политика параметризована вектором $\theta$, а качество политики:
\[
    J(\theta)=\E_{\tau\sim\pi_\theta}R(\tau),
    \qquad
    R(\tau)=\sum_{t=0}^{T}\gamma^t r_t.
\]
Тогда CEM можно применять как black-box оптимизацию параметров политики:
\begin{enumerate}
    \item Задать распределение над параметрами политики $q_\psi(\theta)$.
    \item Насэмплировать $\theta_1,\dots,\theta_N\sim q_\psi$.
    \item Для каждой политики $\pi_{\theta_i}$ сыграть один или несколько эпизодов и оценить $J(\theta_i)$.
    \item Выбрать элитные параметры с наибольшей средней наградой.
    \item Обновить $q_\psi$ по MLE на элитных $\theta_i$.
\end{enumerate}

Если политика линейная или небольшая нейросеть, это может быть вполне рабочим методом. В таком варианте CEM не использует градиенты по политике и не требует знания динамики среды.

\subsection*{CEM для поиска действий или планов}
Другой вариант - применять CEM не к параметрам политики, а к последовательности действий:
\[
    a_{0:H-1}=(a_0,a_1,\dots,a_{H-1}).
\]
Если есть симулятор или модель динамики, можно решать задачу планирования:
\[
    a_{0:H-1}^*\in \argmax_{a_{0:H-1}}\E\left[\sum_{t=0}^{H-1}\gamma^t r(s_t,a_t)\right].
\]
CEM задает распределение над планами, например
\[
    q_\theta(a_{0:H-1}) = \prod_{t=0}^{H-1}\mathcal N(a_t\mid \mu_t,\Sigma_t),
\]
сэмплирует планы, прогоняет их через модель, выбирает элиту и обновляет $\mu_t,\Sigma_t$. На реальном шаге обычно исполняют только первое действие $a_0$, затем получают новое состояние и заново перепланируют. Это вариант MPC (model predictive control).

\subsection*{Связь с MLE и названием метода}
Название ``cross-entropy'' связано с тем, что минимизация
\[
    \KL(p^*\|q_\theta)
\]
эквивалентна минимизации кросс-энтропии между идеальным распределением успешных объектов и модельным распределением:
\[
    H(p^*,q_\theta)=-\E_{p^*}\log q_\theta(X).
\]
А практически это превращается в MLE по элитным сэмплам. Поэтому CEM можно запомнить так: ``сэмплируем - отбираем лучших - обучаем распределение на лучших''.

\subsection*{Что обязательно сказать на экзамене}
Кросс-энтропийный метод - это не один конкретный RL-алгоритм, а общий принцип. Он строит параметрическое распределение над решениями и обновляет его через MLE по элитным сэмплам. В RL решениями могут быть параметры политики, траектории или планы действий. Плюс метода - отсутствие требований к дифференцируемости; минус - высокая сэмпловая стоимость и плохая масштабируемость без дополнительных трюков.

\newpage
\section{Уравнения Беллмана. Policy Iteration и Value Iteration}

\subsection*{MDP и функции ценности}
Марковский процесс принятия решений задается кортежем
\[
    M=(\cS,\cA,p,r,\gamma),
\]
где $\cS$ - состояния, $\cA$ - действия, $p(s'\mid s,a)$ - вероятность перехода, $r(s,a)$ - средняя награда, $\gamma\in(0,1]$ - дисконтирование. Политика $\pi(a\mid s)$ задает распределение действий в состоянии $s$.

Для фиксированной политики $\pi$ вводятся функции ценности:
\[
    V^\pi(s)=\E_\pi\left[\sum_{t=0}^\infty \gamma^t r(s_t,a_t)\mid s_0=s\right],
\]
\[
    Q^\pi(s,a)=\E_\pi\left[\sum_{t=0}^\infty \gamma^t r(s_t,a_t)\mid s_0=s,a_0=a\right].
\]
$V^\pi(s)$ - насколько хорошо начать из состояния $s$ и дальше действовать по $\pi$. $Q^\pi(s,a)$ - насколько хорошо сначала сделать $a$ из $s$, а затем действовать по $\pi$.

\subsection*{Уравнения Беллмана для фиксированной политики}
Связь между $V^\pi$ и $Q^\pi$:
\[
    V^\pi(s)=\sum_{a\in\cA}\pi(a\mid s)Q^\pi(s,a),
\]
\[
    Q^\pi(s,a)=r(s,a)+\gamma\sum_{s'\in\cS}p(s'\mid s,a)V^\pi(s').
\]
Отсюда уравнение Беллмана для $V^\pi$:
\[
    V^\pi(s)=\sum_{a\in\cA}\pi(a\mid s)\left[r(s,a)+\gamma\sum_{s'\in\cS}p(s'\mid s,a)V^\pi(s')\right].
\]
В матричной форме:
\[
    V^\pi = r^\pi + \gamma P^\pi V^\pi,
\]
где
\[
    r^\pi(s)=\sum_a\pi(a\mid s)r(s,a),
    \qquad
    P^\pi(s'\mid s)=\sum_a\pi(a\mid s)p(s'\mid s,a).
\]
Если $\gamma<1$, то
\[
    V^\pi=(I-\gamma P^\pi)^{-1}r^\pi.
\]
На практике часто используют итерации Беллмана:
\[
    V_{k+1}(s)=\sum_a\pi(a\mid s)\left[r(s,a)+\gamma\sum_{s'}p(s'\mid s,a)V_k(s')\right].
\]
Оператор Беллмана для фиксированной политики является сжимающим с коэффициентом $\gamma$ в sup-норме, поэтому итерации сходятся к единственному $V^\pi$.

\subsection*{Оптимальные функции ценности}
Оптимальная ценность:
\[
    V^*(s)=\max_\pi V^\pi(s),
    \qquad
    Q^*(s,a)=\max_\pi Q^\pi(s,a).
\]
Оптимальные уравнения Беллмана:
\[
    V^*(s)=\max_{a\in\cA}\left[r(s,a)+\gamma\sum_{s'}p(s'\mid s,a)V^*(s')\right],
\]
\[
    Q^*(s,a)=r(s,a)+\gamma\sum_{s'}p(s'\mid s,a)\max_{a'\in\cA}Q^*(s',a').
\]
Оптимальная детерминированная политика получается жадно:
\[
    \pi^*(s)\in\argmax_{a\in\cA}Q^*(s,a)
= \argmax_{a\in\cA}\left[r(s,a)+\gamma\sum_{s'}p(s'\mid s,a)V^*(s')\right].
\]
Главная идея: если мы знаем $Q^*$, то задача решена, потому что в каждом состоянии надо выбрать действие с максимальным $Q^*(s,a)$.

\subsection*{Policy Evaluation}
Policy Evaluation - оценивание $V^\pi$ для фиксированной политики $\pi$. Вход: MDP и политика. Выход: $V^\pi$.

Итерационный вариант:
\[
    V_{k+1}=T^\pi V_k,
\]
где
\[
    (T^\pi V)(s)=\sum_a\pi(a\mid s)\left[r(s,a)+\gamma\sum_{s'}p(s'\mid s,a)V(s')\right].
\]
Остановка: например,
\[
    \|V_{k+1}-V_k\|_\infty<\varepsilon.
\]

\subsection*{Policy Improvement}
Policy Improvement - улучшение политики по текущей оценке $V^\pi$ или $Q^\pi$:
\[
    \pi_{\text{new}}(s)\in\argmax_a Q^\pi(s,a).
\]
Если
\[
    Q^\pi(s,\pi_{\text{new}}(s))\ge V^\pi(s)\quad \forall s,
\]
то новая политика не хуже старой:
\[
    V^{\pi_{\text{new}}}(s)\ge V^\pi(s).
\]
Это утверждение называется policy improvement theorem.

\subsection*{Policy Iteration}
Policy Iteration чередует оценивание политики и ее улучшение.

Алгоритм:
\begin{enumerate}
    \item Инициализировать произвольную политику $\pi_0$.
    \item Для $k=0,1,2,\dots$:
    \begin{enumerate}
        \item Policy Evaluation: найти $V^{\pi_k}$ из
        \[
            V^{\pi_k}=T^{\pi_k}V^{\pi_k}.
        \]
        \item Policy Improvement:
        \[
            \pi_{k+1}(s)\in\argmax_a\left[r(s,a)+\gamma\sum_{s'}p(s'\mid s,a)V^{\pi_k}(s')\right].
        \]
        \item Если $\pi_{k+1}=\pi_k$, остановиться: политика оптимальна.
    \end{enumerate}
\end{enumerate}

В конечном MDP с дисконтированием метод сходится за конечное число улучшений, потому что детерминированных политик конечное число, а каждое нетривиальное улучшение делает политику лучше.

\subsection*{Value Iteration}
Value Iteration объединяет оценивание и улучшение в один шаг. Вместо точной оценки $V^\pi$ для каждой политики сразу применяем оптимальный оператор Беллмана:
\[
    V_{k+1}(s)=\max_a\left[r(s,a)+\gamma\sum_{s'}p(s'\mid s,a)V_k(s')\right].
\]
После сходимости извлекаем политику:
\[
    \pi(s)\in\argmax_a\left[r(s,a)+\gamma\sum_{s'}p(s'\mid s,a)V(s')\right].
\]
Оператор
\[
    (T^*V)(s)=\max_a\left[r(s,a)+\gamma\sum_{s'}p(s'\mid s,a)V(s')\right]
\]
тоже является $\gamma$-сжатием, поэтому $V_k\to V^*$.

\subsection*{Сравнение Policy Iteration и Value Iteration}
Policy Iteration обычно делает меньше внешних итераций, но каждая итерация дороже, потому что требует полноценного policy evaluation. Value Iteration делает более дешевые итерации, но их может быть больше. Оба метода требуют знания модели среды: $p(s'\mid s,a)$ и $r(s,a)$, а также конечного или эффективно перебираемого пространства состояний и действий. Поэтому они относятся к динамическому программированию и model-based постановке.

\subsection*{Что обязательно сказать на экзамене}
Уравнения Беллмана - это рекурсивная декомпозиция ``текущая награда плюс дисконтированная будущая ценность''. Для фиксированной политики получаем линейное уравнение $V^\pi=r^\pi+\gamma P^\pi V^\pi$. Для оптимальной политики появляется максимум по действиям. Policy Iteration чередует оценку и улучшение политики; Value Iteration сразу применяет оптимальный оператор Беллмана.

\newpage
\section{Табличные методы: Монте-Карло и Q-learning}

\subsection*{Переход к model-free RL}
В динамическом программировании предполагалось, что известны $p(s'\mid s,a)$ и $r(s,a)$. В model-free RL эти величины неизвестны. Агент может только взаимодействовать со средой и получать переходы
\[
    (s_t,a_t,r_t,s_{t+1}).
\]
Если $|\cS|$ и $|\cA|$ конечны и не слишком велики, можно хранить таблицу значений $V(s)$ или $Q(s,a)$.

\subsection*{Общий вид стохастического обновления}
Многие табличные методы имеют вид экспоненциального сглаживания:
\[
    Q_{k+1}(s,a)=Q_k(s,a)+\alpha_k\left(y_k-Q_k(s,a)\right),
\]
где $y_k$ - целевое значение. Если $y_k$ является шумной оценкой правильного таргета, то такое обновление является стохастической аппроксимацией. Для сходимости обычно нужны условия Роббинса-Монро:
\[
    \sum_{k=1}^\infty \alpha_k=\infty,
    \qquad
    \sum_{k=1}^\infty \alpha_k^2<\infty.
\]
Также нужно достаточно полное исследование среды: каждая пара $(s,a)$ должна посещаться бесконечно часто.

\subsection*{Метод Монте-Карло}
Методы Монте-Карло оценивают ценность по полным эпизодам. Возврат из момента $t$:
\[
    G_t=\sum_{k=t}^{T}\gamma^{k-t}r_k.
\]
Для политики $\pi$:
\[
    Q^\pi(s,a)=\E_\pi[G_t\mid s_t=s,a_t=a].
\]
MC-оценка строится как среднее возвратов после посещения пары $(s,a)$:
\[
    Q(s,a)\approx \frac1{N(s,a)}\sum_{i: (s,a)\text{ visited}}G_i.
\]
Онлайн-обновление:
\[
    Q(s_t,a_t)\leftarrow Q(s_t,a_t)+\alpha\left(G_t-Q(s_t,a_t)\right).
\]

\subsection*{First-visit и every-visit MC}
Есть две стандартные версии:
\begin{itemize}
    \item first-visit MC - в одном эпизоде обновляем $(s,a)$ только по первому посещению;
    \item every-visit MC - обновляем по каждому посещению.
\end{itemize}
Обе версии могут сходиться к $Q^\pi$ при достаточном числе эпизодов.

\subsection*{Плюсы и минусы MC}
Плюсы:
\begin{itemize}
    \item не требует знания модели среды;
    \item дает несмещенную оценку полного возврата;
    \item хорошо распространяет информацию о награде вдоль всей траектории.
\end{itemize}
Минусы:
\begin{itemize}
    \item нужно ждать конца эпизода;
    \item плохо подходит для бесконечных задач;
    \item высокая дисперсия, особенно при длинных эпизодах;
    \item данные используются не очень эффективно.
\end{itemize}

\subsection*{Temporal Difference}
TD-методы не ждут конца эпизода, а используют бутстраппинг:
\[
    V(s_t)\leftarrow V(s_t)+\alpha\left(r_t+\gamma V(s_{t+1})-V(s_t)\right).
\]
TD-ошибка:
\[
    \delta_t=r_t+\gamma V(s_{t+1})-V(s_t).
\]
TD имеет меньшую дисперсию, чем MC, но появляется смещение, потому что таргет использует текущую оценку $V$.

\subsection*{SARSA}
SARSA - on-policy TD-метод для $Q$-функции. Название связано с пятеркой
\[
    (s_t,a_t,r_t,s_{t+1},a_{t+1}).
\]
Обновление:
\[
    Q(s_t,a_t)\leftarrow Q(s_t,a_t)+\alpha\left(r_t+\gamma Q(s_{t+1},a_{t+1})-Q(s_t,a_t)\right),
\]
где $a_{t+1}$ реально выбрано текущей behavior policy. Поэтому SARSA учит ценность той политики, которой агент фактически следует, например $\varepsilon$-жадной политики.

\subsection*{Q-learning}
Q-learning - off-policy TD-метод для обучения оптимальной $Q^*$-функции. Таргет использует жадное действие в следующем состоянии:
\[
    y_t=r_t+\gamma\max_{a'}Q(s_{t+1},a').
\]
Обновление:
\[
    Q(s_t,a_t)\leftarrow Q(s_t,a_t)+\alpha\left(r_t+\gamma\max_{a'}Q(s_{t+1},a')-Q(s_t,a_t)\right).
\]
Соответствующая оптимальная политика:
\[
    \pi(s)\in\argmax_a Q(s,a).
\]

Почему off-policy: действие $a_t$ может быть выбрано исследующей политикой, например $\varepsilon$-greedy, но таргет строится так, как будто дальше агент будет действовать жадно. То есть behavior policy и target policy различаются.

\subsection*{$\varepsilon$-greedy exploration}
Для исследования среды часто используют $\varepsilon$-жадную стратегию:
\[
    a_t=\begin{cases}
    \text{случайное действие}, & \text{с вероятностью }\varepsilon,\\
    \argmax_a Q(s_t,a), & \text{с вероятностью }1-\varepsilon.
    \end{cases}
\]
При убывающем $\varepsilon$ агент сначала исследует, а затем постепенно переходит к эксплуатации найденной политики. Для теоретической сходимости важно, чтобы все пары $(s,a)$ продолжали посещаться достаточно часто.

\subsection*{Сходимость Q-learning}
В табличном случае Q-learning сходится к $Q^*$ при условиях:
\begin{itemize}
    \item конечные $\cS$ и $\cA$;
    \item все пары $(s,a)$ посещаются бесконечно часто;
    \item шаги обучения удовлетворяют условиям Роббинса-Монро;
    \item награды ограничены;
    \item $\gamma<1$ для дисконтированной задачи.
\end{itemize}

\subsection*{MC, SARSA, Q-learning: сравнение}
\begin{center}
\begin{tabular}{>{\raggedright\arraybackslash}p{3cm}p{3.2cm}p{3.2cm}p{3.2cm}}
\toprule
Метод & Таргет & Тип & Особенности \\
\midrule
Monte-Carlo & $G_t$ & обычно on-policy & несмещенный, высокая дисперсия, нужен конец эпизода \\
SARSA & $r+\gamma Q(s',a')$ & on-policy & учитывает реальное исследующее поведение \\
Q-learning & $r+\gamma\max_{a'}Q(s',a')$ & off-policy & учит жадную оптимальную политику \\
\bottomrule
\end{tabular}
\end{center}

\subsection*{Что обязательно сказать на экзамене}
Монте-Карло оценивает $Q^\pi$ по полным возвратам и поэтому имеет низкое смещение, но высокую дисперсию. TD и Q-learning используют бутстраппинг: таргет содержит текущую оценку будущей ценности. Q-learning учит $Q^*$ через таргет $r+\gamma\max_{a'}Q(s',a')$ и является off-policy методом.

\newpage
\section{DQN и модификации: Double DQN, PER, dueling, noisy, multi-step, memory}

\subsection*{Переход от Q-learning к DQN}
Табличный Q-learning неприменим, когда состояний слишком много или они непрерывные, например изображения Atari. Идея DQN: заменить таблицу $Q(s,a)$ нейросетью $Q_\theta(s,a)$.

Если действий мало и они дискретны, удобная архитектура принимает состояние $s$ и выдает вектор:
\[
    Q_\theta(s,\cdot)=\left(Q_\theta(s,a_1),\dots,Q_\theta(s,a_{|\cA|})\right).
\]
Тогда жадное действие:
\[
    a_t=\argmax_a Q_\theta(s_t,a).
\]

\subsection*{Базовый DQN}
Для перехода
\[
    T=(s,a,r,s',done)
\]
строится таргет:
\[
    y(T)=r+\gamma(1-done)\max_{a'}Q_{\theta^-}(s',a'),
\]
где $\theta^-$ - параметры target network. Лосс:
\[
    L(\theta)=\E_{T\sim \cD}\left(Q_\theta(s,a)-y(T)\right)^2.
\]

Алгоритм DQN:
\begin{enumerate}
    \item Инициализировать основную сеть $Q_\theta$, target-сеть $Q_{\theta^-}$ и replay buffer $\cD$.
    \item На каждом шаге выбрать действие через $\varepsilon$-greedy по $Q_\theta$.
    \item Сохранить переход $(s,a,r,s',done)$ в буфер.
    \item Сэмплировать мини-батч переходов из $\cD$.
    \item Построить таргеты $y(T)$ через target network.
    \item Сделать шаг градиентного спуска по MSE/Huber loss.
    \item Периодически обновлять target network: $\theta^-\leftarrow\theta$.
\end{enumerate}

Две ключевые стабилизации DQN:
\begin{itemize}
    \item replay buffer разрушает корреляции между соседними переходами и позволяет переиспользовать опыт;
    \item target network фиксирует таргеты на несколько шагов и уменьшает эффект ``бежим за движущейся целью''.
\end{itemize}

\subsection*{Double DQN}
В vanilla DQN возникает overestimation bias: максимум по шумным оценкам склонен переоценивать истинное значение:
\[
    \E\max_a \hat Q(s,a)\ge \max_a Q(s,a).
\]
Double DQN разделяет выбор действия и оценку действия. Выбор делается основной сетью, оценка - target-сетью:
\[
    a^*(s')=\argmax_{a'}Q_\theta(s',a'),
\]
\[
    y^{DDQN}=r+\gamma(1-done)Q_{\theta^-}(s',a^*(s')).
\]
Идея: одна сеть выбирает argmax, другая оценивает его значение. Это снижает переоценку.

\subsection*{Prioritized Experience Replay}
В обычном replay buffer переходи сэмплируются равномерно. Но не все переходы одинаково полезны. В PER вероятность выбора перехода зависит от TD-ошибки:
\[
    \delta_i = y_i-Q_\theta(s_i,a_i),
\]
\[
    p_i=|\delta_i|+\varepsilon,
    \qquad
    P(i)=\frac{p_i^\alpha}{\sum_j p_j^\alpha}.
\]
Параметр $\alpha$ управляет силой приоритизации: $\alpha=0$ дает равномерный replay.

Такой сэмплинг вносит смещение, поэтому используют importance sampling weights:
\[
    w_i=\left(\frac1{N}\frac1{P(i)}\right)^\beta,
\]
обычно нормируют на максимум в батче. Итоговый лосс:
\[
    L(\theta)=\E_i\left[w_i\left(y_i-Q_\theta(s_i,a_i)\right)^2\right].
\]
Параметр $\beta$ постепенно увеличивают к $1$, чтобы к концу обучения компенсировать смещение сильнее.

\subsection*{Dueling architecture}
Dueling DQN разлагает $Q$ на ценность состояния и преимущество действия:
\[
    Q(s,a)=V(s)+A(s,a).
\]
Но такое разложение неоднозначно: можно прибавить константу к $V$ и вычесть из $A$. Поэтому используют нормировку:
\[
    Q_\theta(s,a)=V_\theta(s)+A_\theta(s,a)-\frac1{|\cA|}\sum_{a'}A_\theta(s,a').
\]
Интуиция: в некоторых состояниях важно знать, что состояние хорошее само по себе, даже если действия почти одинаковы. Dueling-сеть лучше учит $V(s)$ и полезна, когда много действий имеют похожий эффект.

\subsection*{Noisy Networks}
Вместо внешнего $\varepsilon$-greedy можно добавить параметрический шум в веса сети. В noisy linear layer:
\[
    W = \mu_W + \sigma_W\odot\epsilon_W,
    \qquad
    b = \mu_b + \sigma_b\odot\epsilon_b,
\]
где $\epsilon$ - шум, а $\mu,\sigma$ обучаются. Тогда политика случайна из-за случайных параметров сети. Это дает более согласованное исследование: агент не просто случайно дергает действие на каждом шаге, а следует некоторой случайно возмущенной стратегии.

Преимущество над $\varepsilon$-greedy: exploration становится зависящим от состояния и обучаемым.

\subsection*{Multi-step DQN}
Вместо одношагового таргета можно использовать $n$-шаговый возврат:
\[
    y_t^{(n)}=\sum_{k=0}^{n-1}\gamma^k r_{t+k}
    +\gamma^n(1-done)\max_{a'}Q_{\theta^-}(s_{t+n},a').
\]
Плюсы:
\begin{itemize}
    \item награда быстрее распространяется назад;
    \item уменьшается зависимость от неточной bootstrap-оценки на один шаг.
\end{itemize}
Минусы:
\begin{itemize}
    \item выше дисперсия;
    \item сложнее использовать off-policy данные, потому что несколько действий внутри возврата могли быть сделаны старой behavior policy.
\end{itemize}

\subsection*{Память и частичная наблюдаемость}
Обычный DQN предполагает, что состояние $s_t$ марковское. В Atari часто используют frame stacking: состояние составляется из последних кадров, чтобы агент мог восстановить скорость объектов. В общем POMDP текущее наблюдение $o_t$ не содержит всей информации о состоянии. Тогда используют память:
\[
    h_t=f_\theta(h_{t-1},o_t,a_{t-1}),
    \qquad
    Q_\theta(h_t,a).
\]
Практически это DRQN/DQN с LSTM или GRU. Из replay buffer тогда сэмплируют не отдельные переходы, а последовательности, чтобы корректно обучать recurrent state.

\subsection*{Rainbow DQN как комбинация улучшений}
Многие модификации совместимы. Rainbow DQN объединяет несколько идей:
\begin{itemize}
    \item Double DQN;
    \item prioritized replay;
    \item dueling architecture;
    \item multi-step targets;
    \item distributional RL;
    \item noisy nets.
\end{itemize}
Главная мысль: современный DQN - это не один трюк, а набор стабилизаций и улучшений сэмпловой эффективности.

\subsection*{Что обязательно сказать на экзамене}
DQN - это Q-learning с нейросетью, replay buffer и target network. Double DQN борется с overestimation bias. PER чаще показывает сети переходы с большой TD-ошибкой. Dueling architecture отдельно учит ценность состояния и advantage действий. Noisy Nets дают обучаемое исследование. Multi-step targets ускоряют распространение награды. Memory/recurrent DQN нужна для частично наблюдаемых сред.

\newpage
\section{Distributional RL: QR-DQN и Implicit Quantile Networks}

\subsection*{Мотивация}
Обычный RL оценивает математическое ожидание суммарной награды:
\[
    Q^\pi(s,a)=\E[Z^\pi(s,a)],
\]
где случайная величина
\[
    Z^\pi(s,a)=\sum_{t=0}^\infty\gamma^t r_t\mid s_0=s,a_0=a
\]
задает весь reward-to-go. Distributional RL предлагает учить не только среднее, а распределение $Z^\pi(s,a)$.

Причина: разные действия могут иметь одинаковое среднее, но разные риски. Одно действие стабильно дает средний результат, другое иногда дает огромный выигрыш, иногда провал. Обычный $Q$ видит только среднее, а distributional-подход хранит больше информации о неопределенности исходов.

\subsection*{Distributional Bellman Equation}
Для фиксированной политики:
\[
    Z^\pi(s,a) \overset{D}{=} r(s,a)+\gamma Z^\pi(S',A'),
\]
где
\[
    S'\sim p(\cdot\mid s,a),
    \qquad
    A'\sim \pi(\cdot\mid S').
\]
Символ $\overset{D}{=}$ означает равенство по распределению. Это не равенство чисел, а равенство законов случайных величин.

Для оптимального управления:
\[
    Z^*(s,a) \overset{D}{=} r(s,a)+\gamma Z^*(S',a^*),
    \qquad
    a^*\in\argmax_{a'}\E[Z^*(S',a')].
\]
Действие все равно обычно выбирается по среднему:
\[
    a^*(s)=\argmax_a \E[Z(s,a)].
\]

\subsection*{Метрика Вассерштейна}
Для распределений естественно использовать расстояние Вассерштейна. В одномерном случае его удобно выражать через квантильные функции:
\[
    W_p(U,V)=\left(\int_0^1 |F_U^{-1}(\tau)-F_V^{-1}(\tau)|^p\dd\tau\right)^{1/p}.
\]
Distributional Bellman operator является сжимающим в максимальной форме метрики Вассерштейна. Это важная теоретическая мотивация distributional RL.

\subsection*{QR-DQN: представление распределения квантилями}
QR-DQN аппроксимирует распределение $Z(s,a)$ набором $N$ квантильных значений:
\[
    Z_\theta(s,a)=\frac1N\sum_{i=1}^N \delta_{\theta_i(s,a)},
\]
где $\theta_i(s,a)$ - положение $i$-го атома/квантиля. Уровни квантилей фиксированы:
\[
    \tau_i=\frac{i-0.5}{N},\qquad i=1,\dots,N.
\]
Сеть для каждого действия выдает $N$ чисел. Среднее для выбора действия:
\[
    Q(s,a)=\frac1N\sum_{i=1}^N \theta_i(s,a).
\]

\subsection*{Таргет QR-DQN}
Для перехода $(s,a,r,s',done)$ выбираем следующее действие по среднему:
\[
    a^*=\argmax_{a'}\frac1N\sum_{j=1}^N \theta_j^-(s',a').
\]
Таргетные квантили:
\[
    y_j=r+\gamma(1-done)\theta_j^-(s',a^*),
    \qquad j=1,\dots,N.
\]
Теперь надо приблизить набор предсказанных квантилей $\theta_i(s,a)$ к набору таргетных квантилей $y_j$.

\subsection*{Quantile regression loss}
Для квантиля уровня $\tau$ используется pinball loss:
\[
    \rho_\tau(u)=u(\tau-\1\{u<0\}).
\]
Для устойчивости применяют Huber-версию:
\[
    \rho_\tau^\kappa(u)=|\tau-\1\{u<0\}|\mathcal L_\kappa(u),
\]
где
\[
    \mathcal L_\kappa(u)=
    \begin{cases}
    \frac12u^2, & |u|\le \kappa,\\
    \kappa(|u|-\frac12\kappa), & |u|>\kappa.
    \end{cases}
\]
Ошибка между предсказанными и таргетными квантилями:
\[
    L(\theta)=\frac1N\sum_{i=1}^N\sum_{j=1}^N
    \rho_{\tau_i}^\kappa\left(y_j-\theta_i(s,a)\right).
\]

\subsection*{Алгоритм QR-DQN}
\begin{enumerate}
    \item Сеть выдает квантили $\theta_i(s,a)$ для каждого действия.
    \item Действие выбирается $\varepsilon$-жадно по среднему квантилей.
    \item Переход сохраняется в replay buffer.
    \item Для батча переходов строятся distributional targets.
    \item Сеть обучается quantile regression loss.
    \item Target network обновляется как в DQN.
\end{enumerate}

QR-DQN отличается от DQN не тем, как выбирается действие, а тем, что вместо одного числа $Q(s,a)$ учится распределение будущего возврата.

\subsection*{Implicit Quantile Networks}
В QR-DQN уровни квантилей фиксированы. IQN делает сильнее: сеть принимает не только состояние, но и уровень квантиля $\tau\in(0,1)$:
\[
    Z_\theta(s,a;\tau)\approx F^{-1}_{Z(s,a)}(\tau).
\]
То есть IQN аппроксимирует всю квантильную функцию, а не только фиксированный набор точек.

Обычно $\tau$ кодируется признаками, например через cosine embedding:
\[
    \phi(\tau)=\mathrm{ReLU}\left(W[\cos(\pi i\tau)]_{i=1}^n+b\right),
\]
после чего embedding состояния перемножается или объединяется с embedding квантиля. Сеть выдает значение квантиля для каждого действия.

\subsection*{Таргет и лосс IQN}
Сэмплируем уровни квантилей $\tau_i$ для текущей сети и $\tau'_j$ для таргета:
\[
    \tau_i,\tau'_j\sim U(0,1).
\]
Выбор действия:
\[
    a^*=\argmax_{a'}\frac1{N'}\sum_{j=1}^{N'}Z_{\theta^-}(s',a';\tau'_j).
\]
Таргеты:
\[
    y_j=r+\gamma(1-done)Z_{\theta^-}(s',a^*;\tau'_j).
\]
Лосс:
\[
    L(\theta)=\frac1{NN'}\sum_{i=1}^N\sum_{j=1}^{N' }
    \rho_{\tau_i}^\kappa\left(y_j-Z_\theta(s,a;\tau_i)\right).
\]

\subsection*{Почему IQN лучше QR-DQN}
QR-DQN хранит фиксированную дискретизацию распределения. IQN может сэмплировать любые уровни квантилей и тем самым аппроксимирует распределение гибче. Это дает более точную оценку хвостов распределения и лучше масштабируется по числу квантильных точек: можно менять количество сэмплов $\tau$ на обучении и инференсе.

\subsection*{Risk-sensitive политики}
Distributional RL позволяет выбирать действия не только по среднему. Например:
\begin{itemize}
    \item risk-averse: оптимизировать нижние квантили;
    \item risk-seeking: оптимизировать верхние квантили;
    \item CVaR: среднее по худшему хвосту распределения.
\end{itemize}
Но в базовых QR-DQN/IQN действие обычно выбирают по математическому ожиданию.

\subsection*{Что обязательно сказать на экзамене}
Distributional RL учит распределение возврата $Z(s,a)$, а не только его среднее $Q(s,a)$. Основное уравнение: $Z(s,a)\overset{D}{=}r+\gamma Z(s',a')$. QR-DQN представляет распределение фиксированным набором квантилей и обучает их quantile regression loss. IQN принимает уровень квантиля $\tau$ на вход и аппроксимирует всю квантильную функцию.

\newpage
\section{Внутренняя мотивация: RND и ICM}

\subsection*{Проблема exploration}
В задачах с разреженной наградой агент долго не получает полезный сигнал. Например, награда $+1$ выдается только при достижении цели, а во всех остальных состояниях $0$ или маленький штраф. Тогда случайное исследование может почти никогда не привести к цели.

Идея внутренней мотивации: добавить к внешней награде $r^{extr}$ дополнительную внутреннюю награду $r^{intr}$, которая поощряет исследование:
\[
    r(s,a)=r^{extr}(s,a)+\alpha r^{intr}(s,a),
\]
где $\alpha$ - коэффициент масштаба.

Для фиксированной политики суммы ценностей раскладываются:
\[
    Q^\pi(s,a)=Q^\pi_{extr}(s,a)+\alpha Q^\pi_{intr}(s,a),
\]
\[
    V^\pi(s)=V^\pi_{extr}(s)+\alpha V^\pi_{intr}(s).
\]
Но для оптимальных $Q^*,V^*$ такое разложение в общем случае неверно, потому что максимум суммы не равен сумме максимумов.

\subsection*{Простейшие count-based бонусы}
В табличных средах можно поощрять новые состояния:
\[
    r^{intr}(s)=\frac{1}{\sqrt{N(s)}}
\]
или похожей убывающей функцией от числа посещений $N(s)$. Чем реже состояние посещалось, тем больше бонус.

В больших пространствах состояний точные счетчики не работают. Поэтому используют приближения: хэширование, density models, prediction error. RND и ICM относятся к prediction-based подходам.

\subsection*{Random Network Distillation}
RND оценивает новизну состояния через ошибку предсказания случайной фиксированной сети.

Есть две сети:
\begin{itemize}
    \item target network $f(s)$ - случайно инициализированная и замороженная;
    \item predictor network $\hat f_\theta(s)$ - обучаемая сеть, которая пытается предсказывать выход target network.
\end{itemize}

Внутренняя награда:
\[
    r_t^{intr}=\left\|\hat f_\theta(s_t)-f(s_t)\right\|^2.
\]
Если состояние часто встречалось, predictor уже научился предсказывать $f(s)$, ошибка мала. Если состояние новое, ошибка велика, агент получает бонус.

Лосс обучения predictor:
\[
    L_{RND}(\theta)=\E_{s\sim\cD}\left\|\hat f_\theta(s)-f(s)\right\|^2.
\]

\subsection*{Почему target network случайная}
Случайная сеть $f$ задает фиксированное случайное embedding-пространство. Агент не должен оптимизировать саму задачу предсказания target network через изменение target network: цель должна быть неподвижной. Поэтому обучается только predictor.

Такой подход похож на дистилляцию: predictor пытается воспроизвести выходы случайного учителя.

\subsection*{Преимущества и проблемы RND}
Преимущества:
\begin{itemize}
    \item простая реализация;
    \item работает в высокоразмерных наблюдениях;
    \item не требует модели действий;
    \item хорошо подходит для sparse reward.
\end{itemize}

Проблемы:
\begin{itemize}
    \item ошибка предсказания зависит от масштаба признаков, нужна нормализация;
    \item внутреннюю награду надо нормировать;
    \item стохастический шум в наблюдениях может выглядеть как новизна;
    \item при долгом обучении бонусы исчезают, и агент может перестать исследовать.
\end{itemize}

\subsection*{Intrinsic Curiosity Module}
ICM строит внутреннюю награду через ошибку предсказания следующего состояния в learned feature space. Основная идея: агенту интересно то, что он плохо предсказывает, но что зависит от его действий.

ICM содержит:
\begin{itemize}
    \item энкодер $\phi(s)$, переводящий состояние в признаки;
    \item inverse dynamics model: по $\phi(s_t)$ и $\phi(s_{t+1})$ предсказывает действие $a_t$;
    \item forward dynamics model: по $\phi(s_t)$ и $a_t$ предсказывает $\phi(s_{t+1})$.
\end{itemize}

Inverse model:
\[
    \hat a_t = g_\eta(\phi(s_t),\phi(s_{t+1})).
\]
Лосс inverse model для дискретных действий обычно cross-entropy:
\[
    L_{inv}= -\log p_\eta(a_t\mid \phi(s_t),\phi(s_{t+1})).
\]

Forward model:
\[
    \widehat{\phi(s_{t+1})}=F_\psi(\phi(s_t),a_t).
\]
Лосс forward model:
\[
    L_{fwd}=\frac12\left\|F_\psi(\phi(s_t),a_t)-\phi(s_{t+1})\right\|^2.
\]

Внутренняя награда:
\[
    r_t^{intr}=\frac\eta2\left\|F_\psi(\phi(s_t),a_t)-\phi(s_{t+1})\right\|^2.
\]

\subsection*{Зачем нужна inverse model}
Если просто предсказывать следующий кадр, агент может получать большую внутреннюю награду за неуправляемый шум: мерцание телевизора, случайные текстуры, другие агенты. Inverse model заставляет энкодер $\phi$ сохранять информацию, связанную с действием агента. Если по двум соседним признакам нельзя восстановить действие, такие признаки не очень полезны для управляемой динамики.

Иными словами, inverse model помогает игнорировать неуправляемые аспекты наблюдения, а forward model измеряет ошибку предсказания в пространстве управляемых признаков.

\subsection*{Общий лосс ICM}
Обычно оптимизируют комбинацию:
\[
    L_{ICM}=(1-\beta)L_{inv}+\beta L_{fwd},
\]
а агент максимизирует награду
\[
    r_t=r_t^{extr}+\alpha r_t^{intr}.
\]
Параметр $\beta$ балансирует обучение признаков через inverse и forward задачи.

\subsection*{RND против ICM}
RND измеряет новизну состояния: ``я раньше не видел такое состояние''. ICM измеряет непредсказуемость перехода, зависящего от действия: ``я не понимаю, что произойдет после моего действия''. RND проще, но хуже отделяет управляемую новизну от случайного шума. ICM сложнее, но через inverse dynamics пытается фокусироваться на аспектах среды, контролируемых агентом.

\subsection*{Что обязательно сказать на экзамене}
Внутренняя мотивация добавляет к внешней награде бонус за исследование. RND дает бонус как ошибку предсказания случайной замороженной сети. ICM дает бонус как ошибку forward model в feature space, а inverse model обучает признаки быть связанными с действиями агента. Главная цель - помочь exploration в sparse reward задачах.

\newpage
\section{Policy Gradient. REINFORCE и модификации}

\subsection*{Зачем нужны policy gradient методы}
Value-based методы учат $Q(s,a)$ и выбирают действие через $\argmax_a Q(s,a)$. Это удобно для дискретных действий, но сложно для непрерывных: максимум по действиям может быть дорогим. Policy gradient методы напрямую параметризуют политику:
\[
    \pi_\theta(a\mid s)
\]
и оптимизируют ожидаемую награду:
\[
    J(\theta)=\E_{\tau\sim\pi_\theta}\left[\sum_{t=0}^\infty\gamma^t r_t\right].
\]

\subsection*{Likelihood-ratio trick}
Пусть вероятность траектории:
\[
    p_\theta(\tau)=p(s_0)\prod_{t\ge0}\pi_\theta(a_t\mid s_t)p(s_{t+1}\mid s_t,a_t).
\]
Динамика среды не зависит от $\theta$, поэтому
\[
    \nabla_\theta \log p_\theta(\tau)=\sum_{t\ge0}\nabla_\theta\log \pi_\theta(a_t\mid s_t).
\]
Используем трюк:
\[
    \nabla_\theta p_\theta(\tau)=p_\theta(\tau)\nabla_\theta\log p_\theta(\tau).
\]
Тогда
\[
    \nabla_\theta J(\theta)=\E_{\tau\sim\pi_\theta}\left[\nabla_\theta\log p_\theta(\tau)R(\tau)\right].
\]

\subsection*{Policy Gradient Theorem}
Классическая форма:
\[
    \nabla_\theta J(\theta)=\frac1{1-\gamma}\E_{s\sim d^{\pi_\theta},a\sim\pi_\theta}\left[\nabla_\theta\log\pi_\theta(a\mid s)Q^{\pi_\theta}(s,a)\right],
\]
где дисконтированное распределение посещения состояний:
\[
    d^\pi(s)=(1-\gamma)\sum_{t=0}^\infty\gamma^t\Prob(s_t=s\mid\pi).
\]
Эквивалентная траекторная форма:
\[
    \nabla_\theta J(\theta)=\E_{\tau\sim\pi_\theta}\sum_{t\ge0}\gamma^t
    \nabla_\theta\log\pi_\theta(a_t\mid s_t)Q^{\pi_\theta}(s_t,a_t).
\]

\subsection*{REINFORCE}
В REINFORCE неизвестное $Q^\pi(s_t,a_t)$ заменяется Монте-Карло возвратом:
\[
    G_t=\sum_{k=t}^{T}\gamma^{k-t}r_k.
\]
Оценка градиента:
\[
    \widehat{\nabla_\theta J}=\frac1N\sum_{i=1}^N\sum_{t=0}^{T_i}\gamma^t
    \nabla_\theta\log\pi_\theta(a_t^i\mid s_t^i)G_t^i.
\]
Параметры обновляют градиентным подъемом:
\[
    \theta\leftarrow\theta+\alpha\widehat{\nabla_\theta J}.
\]
В терминах loss для градиентного спуска часто пишут:
\[
    L(\theta)=-\sum_t \log\pi_\theta(a_t\mid s_t)G_t.
\]

\subsection*{Интуиция REINFORCE}
Если действие привело к большому будущему возврату $G_t$, то увеличиваем $\log\pi_\theta(a_t\mid s_t)$, то есть повышаем вероятность этого действия. Если возврат малый или отрицательный, вероятность действия уменьшается. Это прямой метод ``усиливать действия, которые встретились в хороших траекториях''.

\subsection*{Проблема дисперсии}
REINFORCE имеет несмещенную, но очень шумную оценку градиента:
\begin{itemize}
    \item нужно ждать завершения эпизода;
    \item возврат $G_t$ может иметь большую дисперсию;
    \item одно действие получает кредит за весь будущий эпизод, включая случайность среды.
\end{itemize}
Поэтому чистый REINFORCE редко используется в сложных задачах без модификаций.

\subsection*{Baseline}
Можно вычитать из $Q$ любую функцию состояния $b(s)$:
\[
    \nabla_\theta J(\theta)=\frac1{1-\gamma}\E_{s,a}\left[\nabla_\theta\log\pi_\theta(a\mid s)(Q^\pi(s,a)-b(s))\right].
\]
Это не меняет матожидание градиента, потому что
\[
    \E_{a\sim\pi}\nabla_\theta\log\pi_\theta(a\mid s)b(s)=b(s)\nabla_\theta\sum_a\pi_\theta(a\mid s)=0.
\]
На практике выбирают
\[
    b(s)=V^\pi(s).
\]
Тогда появляется advantage:
\[
    A^\pi(s,a)=Q^\pi(s,a)-V^\pi(s).
\]

\subsection*{REINFORCE with baseline}
Оценка градиента:
\[
    \widehat{\nabla_\theta J}=\sum_t\gamma^t
    \nabla_\theta\log\pi_\theta(a_t\mid s_t)(G_t-b(s_t)).
\]
Если $b(s)=V_\phi(s)$, то критик обучается регрессией:
\[
    L_V(\phi)=\sum_t\left(V_\phi(s_t)-G_t\right)^2.
\]

\subsection*{Reward-to-go вместо полного возврата}
Вместо полного возврата всей траектории $R(\tau)$ для каждого действия используют только будущий возврат $G_t$. Это снижает дисперсию: действие $a_t$ не должно получать кредит за награды, случившиеся до него.

\subsection*{Entropy regularization}
Чтобы политика не схлопывалась слишком рано, добавляют бонус энтропии:
\[
    J_{ent}(\theta)=J(\theta)+\beta\E_{s\sim d^\pi}H(\pi_\theta(\cdot\mid s)),
\]
где
\[
    H(\pi(\cdot\mid s))=-\E_{a\sim\pi}\log\pi(a\mid s).
\]
Это поощряет exploration и делает распределение действий менее вырожденным.

\subsection*{On-policy характер}
Policy gradient в базовом виде on-policy: данные должны быть собраны текущей политикой $\pi_\theta$. Если данные собраны старой политикой, оценка градиента для новой политики становится смещенной. Позднее это исправляют через importance sampling, TRPO/PPO и actor-critic схемы.

\subsection*{Что обязательно сказать на экзамене}
Policy Gradient напрямую оптимизирует параметры политики. Основная формула: $\nabla J=\E[\nabla\log\pi(a\mid s)Q(s,a)]$. REINFORCE заменяет $Q$ на MC-return $G_t$. Главная проблема - высокая дисперсия. Baseline $V(s)$ не меняет матожидание градиента, но снижает дисперсию, приводя к advantage $A=Q-V$.

\newpage
\section{Actor-Critic. A2C и GAE}

\subsection*{Идея Actor-Critic}
Actor-Critic объединяет policy gradient и value learning. Есть две модели:
\begin{itemize}
    \item actor $\pi_\theta(a\mid s)$ - политика, выбирающая действия;
    \item critic $V_\phi(s)$ или $Q_\phi(s,a)$ - оценочная функция, оценивающая действия или состояния.
\end{itemize}
Actor обновляется по policy gradient, critic обучается предсказывать ценность и используется для уменьшения дисперсии.

Базовая формула градиента:
\[
    \nabla_\theta J(\theta)=\E\left[\nabla_\theta\log\pi_\theta(a_t\mid s_t)A^\pi(s_t,a_t)\right].
\]
Actor-Critic заменяет истинный $A^\pi$ оценкой критика.

\subsection*{TD-оценка advantage}
Если critic оценивает $V_\phi(s)$, то одношаговая TD-ошибка:
\[
    \delta_t=r_t+\gamma V_\phi(s_{t+1})-V_\phi(s_t).
\]
Она является оценкой advantage:
\[
    \delta_t\approx A^\pi(s_t,a_t).
\]
Тогда actor update:
\[
    \nabla_\theta J\approx \sum_t \nabla_\theta\log\pi_\theta(a_t\mid s_t)\delta_t.
\]
Critic обучается по таргету
\[
    y_t=r_t+\gamma V_\phi(s_{t+1}),
\]
лосс:
\[
    L_V(\phi)=\sum_t\left(V_\phi(s_t)-y_t\right)^2.
\]
На практике в таргете градиент через $V_\phi(s_{t+1})$ останавливают.

\subsection*{A2C}
A2C - Advantage Actor-Critic. Это синхронная версия A3C: несколько параллельных окружений собирают rollout, затем один общий батч используется для обновления actor и critic.

Общий loss часто записывают так:
\[
    L(\theta,\phi)=L_{policy}+c_vL_{value}-c_eL_{entropy},
\]
где
\[
    L_{policy}(\theta)=-\sum_t \log\pi_\theta(a_t\mid s_t)\hat A_t,
\]
\[
    L_{value}(\phi)=\sum_t\left(V_\phi(s_t)-\hat R_t\right)^2,
\]
\[
    L_{entropy}=\sum_t H(\pi_\theta(\cdot\mid s_t)).
\]
Знак перед entropy отрицательный в loss, потому что при минимизации мы хотим максимизировать энтропию.

\subsection*{$n$-step returns в A2C}
Вместо одношагового TD можно использовать $n$-шаговый таргет:
\[
    \hat R_t^{(n)}=\sum_{l=0}^{n-1}\gamma^l r_{t+l}+\gamma^n V_\phi(s_{t+n}).
\]
Advantage:
\[
    \hat A_t^{(n)}=\hat R_t^{(n)}-V_\phi(s_t).
\]
При $n=1$ получаем TD(0), при $n\to\infty$ - Monte-Carlo return без bootstrap. Это уже показывает bias-variance trade-off.

\subsection*{Bias-variance trade-off}
Одношаговый TD:
\begin{itemize}
    \item низкая дисперсия;
    \item высокий bias, потому что сильно зависит от текущего критика.
\end{itemize}
Monte-Carlo:
\begin{itemize}
    \item низкий bias;
    \item высокая дисперсия.
\end{itemize}
$n$-step return находится между ними. Нужно выбрать $n$, но фиксированный $n$ не всегда оптимален.

\subsection*{GAE: Generalized Advantage Estimation}
GAE объединяет разные $n$-step оценки advantage с экспоненциальными весами. Сначала вводим TD residual:
\[
    \delta_t^V=r_t+\gamma V(s_{t+1})-V(s_t).
\]
GAE:
\[
    \hat A_t^{GAE(\gamma,\lambda)}=
    \sum_{l=0}^{\infty}(\gamma\lambda)^l\delta_{t+l}^V.
\]
Параметр $\lambda\in[0,1]$ управляет компромиссом:
\begin{itemize}
    \item $\lambda=0$: одношаговый TD advantage $\hat A_t=\delta_t$;
    \item $\lambda=1$: оценка близка к Monte-Carlo advantage;
    \item промежуточные значения дают баланс bias и variance.
\end{itemize}

\subsection*{Связь GAE с $n$-step advantage}
$n$-step advantage:
\[
    \hat A_t^{(n)}=\sum_{l=0}^{n-1}\gamma^l r_{t+l}+\gamma^n V(s_{t+n})-V(s_t).
\]
Его можно представить как сумму TD-ошибок:
\[
    \hat A_t^{(n)}=\sum_{l=0}^{n-1}\gamma^l\delta_{t+l}^V.
\]
GAE - это экспоненциально взвешенная смесь таких оценок:
\[
    \hat A_t^{GAE}=(1-\lambda)\left(\hat A_t^{(1)}+\lambda\hat A_t^{(2)}+\lambda^2\hat A_t^{(3)}+\dots\right).
\]
Эквивалентная компактная форма - сумма TD residuals с весами $(\gamma\lambda)^l$.

\subsection*{Почему Actor-Critic может быть on-policy и off-policy}
Классический policy gradient требует on-policy данные. Но critic часто можно обучать off-policy, потому что Bellman targets используют переходы из буфера. Поэтому actor-critic - это семейство методов. A2C/PPO обычно on-policy. DDPG/TD3/SAC - off-policy actor-critic: critic учится по replay buffer, actor улучшается через critic.

\subsection*{Преимущества Actor-Critic}
По сравнению с REINFORCE:
\begin{itemize}
    \item ниже дисперсия за счет критика;
    \item можно обновляться до конца эпизода;
    \item лучше сэмпловая эффективность;
    \item естественная работа с непрерывными действиями через actor.
\end{itemize}
Недостаток: появляется bias из-за неточного критика. Если critic плохой, actor получает неправильный сигнал.

\subsection*{Что обязательно сказать на экзамене}
Actor-Critic использует actor для политики и critic для оценки advantage. A2C оптимизирует policy loss, value loss и entropy bonus. Advantage можно оценивать через TD residual, $n$-step return или GAE. GAE задается формулой $\hat A_t=\sum_l(\gamma\lambda)^l\delta_{t+l}$ и управляет bias-variance компромиссом через $\lambda$.

\newpage
\section{TRPO: Trust Region Policy Optimization}

\subsection*{Проблема больших шагов policy gradient}
Обычный policy gradient делает шаг по параметрам политики:
\[
    \theta_{k+1}=\theta_k+\alpha\widehat{\nabla J(\theta_k)}.
\]
Но маленькое изменение параметров не всегда означает маленькое изменение политики. Если политика изменилась слишком сильно, собранные старой политикой данные становятся нерелевантны, а качество может резко упасть. TRPO решает это через ограничение на изменение политики.

\subsection*{Relative Performance Identity}
Для двух политик $\pi$ и $\pi_{old}$ выполняется связь:
\[
    J(\pi)-J(\pi_{old})=\frac1{1-\gamma}\E_{s\sim d^\pi,a\sim\pi}\left[A^{\pi_{old}}(s,a)\right],
\]
где
\[
    A^{\pi_{old}}(s,a)=Q^{\pi_{old}}(s,a)-V^{\pi_{old}}(s).
\]
Это точная формула, но неудобная: матожидание берется по $d^\pi$ - распределению состояний новой политики, которого у нас еще нет.

\subsection*{Суррогатный функционал}
TRPO заменяет распределение состояний новой политики на старое:
\[
    L_{\pi_{old}}(\pi)=J(\pi_{old})+\frac1{1-\gamma}
    \E_{s\sim d^{\pi_{old}},a\sim\pi}\left[A^{\pi_{old}}(s,a)\right].
\]
Так как действия в данных собраны из $\pi_{old}$, применяют importance sampling:
\[
    \E_{a\sim\pi}[A(s,a)]
    =\E_{a\sim\pi_{old}}\left[\frac{\pi(a\mid s)}{\pi_{old}(a\mid s)}A(s,a)\right].
\]
Итого практический surrogate:
\[
    L_{\pi_{old}}(\theta)=\E_{s,a\sim\pi_{old}}
    \left[\frac{\pi_\theta(a\mid s)}{\pi_{old}(a\mid s)}A^{\pi_{old}}(s,a)\right].
\]
Обозначение ratio:
\[
    \rho_\theta(s,a)=\frac{\pi_\theta(a\mid s)}{\pi_{old}(a\mid s)}.
\]

\subsection*{Trust region}
Чтобы суррогат был надежным, новая политика не должна сильно отличаться от старой. TRPO решает задачу:
\[
    \max_\theta \; \E_{s,a\sim\pi_{old}}
    \left[\rho_\theta(s,a)A^{\pi_{old}}(s,a)\right]
\]
при ограничении
\[
    \E_{s\sim d^{\pi_{old}}}\left[\KL\left(\pi_{old}(\cdot\mid s)\|\pi_\theta(\cdot\mid s)\right)\right]
    \le \delta.
\]
$\delta$ - радиус trust region. Он ограничивает среднее KL-расстояние между старой и новой политикой.

\subsection*{Теоретическая гарантия}
Основная идея теории TRPO: можно получить нижнюю оценку на качество новой политики через суррогатный функционал и штраф за отличие политик. В грубом виде:
\[
    J(\pi)\ge L_{\pi_{old}}(\pi)-C\cdot D_{KL}^{max}(\pi_{old},\pi),
\]
где $C$ зависит от максимального advantage и $\gamma$. Поэтому если максимизировать surrogate и не уходить далеко по KL, можно ожидать монотонного улучшения.

Практически TRPO заменяет сложное ограничение на максимальный KL средним KL и решает приближенную constrained optimization задачу.

\subsection*{Квадратичная аппроксимация ограничения}
Пусть $g=\nabla_\theta L(\theta)|_{\theta=\theta_{old}}$. KL около старых параметров аппроксимируется через матрицу Фишера $F$:
\[
    \KL(\theta_{old},\theta_{old}+x)\approx \frac12x^\top F x.
\]
Тогда локальная задача:
\[
    \max_x g^\top x
    \quad\text{s.t.}\quad
    \frac12x^\top F x\le\delta.
\]
Решение направлено вдоль natural gradient:
\[
    x^* \propto F^{-1}g.
\]
С учетом ограничения:
\[
    x^*=\sqrt{\frac{2\delta}{g^\top F^{-1}g}}F^{-1}g.
\]

\subsection*{Natural gradient}
Обычный градиент измеряет шаг в евклидовом пространстве параметров. Natural gradient учитывает геометрию распределений политик:
\[
    \tilde\nabla J = F^{-1}\nabla J.
\]
Он делает шаг, который более осмыслен в пространстве вероятностных распределений, а не в пространстве параметров нейросети.

\subsection*{Практический алгоритм TRPO}
\begin{enumerate}
    \item Собрать trajectories текущей политикой $\pi_{old}$.
    \item Оценить $A^{\pi_{old}}(s,a)$, например через GAE.
    \item Построить surrogate objective
    \[
        L(\theta)=\E\left[\rho_\theta(s,a)\hat A(s,a)\right].
    \]
    \item Посчитать градиент $g=\nabla_\theta L$.
    \item Приближенно решить $Fx=g$ через conjugate gradient, не формируя $F$ явно.
    \item Выбрать длину шага, удовлетворяющую KL-ограничению.
    \item Сделать line search: принять шаг, если surrogate вырос и KL не превышает $\delta$.
    \item Обновить critic отдельно.
\end{enumerate}

\subsection*{Плюсы и минусы TRPO}
Плюсы:
\begin{itemize}
    \item стабильнее обычного policy gradient;
    \item теоретически мотивирует монотонное улучшение;
    \item контролирует размер обновления политики.
\end{itemize}
Минусы:
\begin{itemize}
    \item сложная реализация;
    \item нужен conjugate gradient и line search;
    \item метод on-policy и требует много новых данных;
    \item дороже, чем PPO.
\end{itemize}

\subsection*{Что обязательно сказать на экзамене}
TRPO максимизирует surrogate $\E[\frac{\pi_\theta}{\pi_{old}}A_{old}]$, но ограничивает средний KL между новой и старой политикой. Это trust region: новая политика не должна уходить далеко, чтобы surrogate оставался надежной оценкой улучшения. Локально TRPO приводит к natural gradient с ограничением $\frac12x^TFx\le\delta$.

\newpage
\section{Bias-variance, GAE и PPO}

\subsection*{Bias-variance trade-off в RL}
В actor-critic методах нужно оценивать advantage. Есть два крайних варианта.

Monte-Carlo:
\[
    \hat A_t^{MC}=\sum_{l=0}^{T-t}\gamma^l r_{t+l}-V(s_t).
\]
Он почти не имеет bias, если траектория полная и $V$ не используется в таргете, но имеет большую дисперсию.

TD(0):
\[
    \hat A_t^{TD}=r_t+\gamma V(s_{t+1})-V(s_t)=\delta_t.
\]
Он имеет меньшую дисперсию, но больший bias из-за bootstrap через неточный $V$.

$n$-step оценка:
\[
    \hat A_t^{(n)}=\sum_{l=0}^{n-1}\gamma^l r_{t+l}+\gamma^nV(s_{t+n})-V(s_t).
\]
Чем больше $n$, тем меньше bias и больше variance. Чем меньше $n$, тем больше зависимость от критика и ниже variance.

\subsection*{GAE}
Generalized Advantage Estimation вводит TD residual:
\[
    \delta_t=r_t+\gamma V(s_{t+1})-V(s_t).
\]
И определяет:
\[
    \hat A_t^{GAE(\gamma,\lambda)}=\sum_{l=0}^{\infty}(\gamma\lambda)^l\delta_{t+l}.
\]
Параметр $\lambda$:
\begin{itemize}
    \item $\lambda=0$: только одношаговая TD-ошибка;
    \item $\lambda=1$: почти Monte-Carlo advantage;
    \item $0<\lambda<1$: компромисс.
\end{itemize}
На практике часто используют $\lambda\approx0.95$.

Возврат для обучения critic:
\[
    \hat R_t=\hat A_t+V(s_t).
\]

\subsection*{Переход от TRPO к PPO}
TRPO стабилен, но сложен. PPO сохраняет идею ``не менять политику слишком сильно'', но заменяет constrained optimization на простой loss, который можно оптимизировать обычным SGD/Adam.

Ratio новой и старой политики:
\[
    \rho_t(\theta)=\frac{\pi_\theta(a_t\mid s_t)}{\pi_{old}(a_t\mid s_t)}.
\]
Суррогат без ограничения:
\[
    L^{PG}(\theta)=\E_t[\rho_t(\theta)\hat A_t].
\]
Если $\rho_t$ сильно отличается от $1$, политика слишком сильно изменила вероятность действия.

\subsection*{PPO с клиппированием}
Основной clipped objective:
\[
    L^{CLIP}(\theta)=\E_t\left[
    \min\left(
    \rho_t(\theta)\hat A_t,
    \clip(\rho_t(\theta),1-\epsilon,1+\epsilon)\hat A_t
    \right)
    \right].
\]
Здесь $\epsilon$ обычно $0.1$--$0.3$.

Интуиция по знаку advantage:
\begin{itemize}
    \item Если $\hat A_t>0$, действие было лучше ожидаемого. Мы хотим увеличить его вероятность, но не более чем в $1+\epsilon$ раз.
    \item Если $\hat A_t<0$, действие было хуже ожидаемого. Мы хотим уменьшить его вероятность, но не слишком сильно: не ниже $1-\epsilon$.
\end{itemize}
Клиппирование обнуляет стимул делать слишком большой шаг в направлении, которое уже достаточно изменило политику.

\subsection*{Полный PPO loss}
Обычно PPO оптимизирует:
\[
    L(\theta,\phi)= -L^{CLIP}(\theta)+c_1L^{VF}(\phi)-c_2L^{ENT}(\theta),
\]
где
\[
    L^{VF}(\phi)=\E_t\left(V_\phi(s_t)-\hat R_t\right)^2,
\]
\[
    L^{ENT}(\theta)=\E_t\left[H(\pi_\theta(\cdot\mid s_t))\right].
\]
При минимизации знак перед $L^{CLIP}$ отрицательный, потому что его надо максимизировать.

\subsection*{Клиппирование критика}
В некоторых реализациях клиппируют value update:
\[
    V_{\phi}^{clip}(s_t)=V_{old}(s_t)+\clip(V_\phi(s_t)-V_{old}(s_t),-\epsilon_v,\epsilon_v).
\]
Value loss:
\[
    L^{VF}=\max\left(
    (V_\phi(s_t)-\hat R_t)^2,
    (V_\phi^{clip}(s_t)-\hat R_t)^2
    \right).
\]
Цель та же: не позволять критику слишком резко менять оценки.

\subsection*{PPO с KL penalty}
Есть альтернативная версия PPO:
\[
    L^{KLPEN}(\theta)=\E_t[\rho_t(\theta)\hat A_t]
    -\beta\E_t\left[\KL(\pi_{old}(\cdot\mid s_t)\|\pi_\theta(\cdot\mid s_t))\right].
\]
Можно адаптировать $\beta$: если KL слишком большой, увеличить штраф; если слишком маленький, уменьшить. Но clipped PPO обычно проще и популярнее.

\subsection*{Алгоритм PPO}
\begin{enumerate}
    \item Собрать rollout текущей политикой $\pi_{old}$.
    \item Оценить значения $V(s_t)$ и advantage $\hat A_t$, обычно через GAE.
    \item Посчитать returns $\hat R_t=\hat A_t+V(s_t)$.
    \item Зафиксировать log-prob старой политики $\log\pi_{old}(a_t\mid s_t)$.
    \item Несколько эпох оптимизировать clipped objective на этом же батче.
    \item Обновить $\pi_{old}\leftarrow\pi_\theta$ и собрать новый rollout.
\end{enumerate}

PPO остается on-policy: один rollout можно переиспользовать несколько эпох, но нельзя слишком долго обучаться на старых данных, потому что ratio и KL станут ненадежными.

\subsection*{Почему PPO стал популярным}
PPO проще TRPO: не нужны conjugate gradient и line search. При этом он сохраняет ключевую идею ограничения шага политики. Метод достаточно стабилен, хорошо работает в дискретных и непрерывных действиях, легко реализуется и стал стандартным baseline для on-policy deep RL.

\subsection*{Что обязательно сказать на экзамене}
Bias-variance trade-off возникает при оценке advantage: MC имеет низкий bias и высокую variance, TD - наоборот. GAE задает сглаженную сумму TD-ошибок $\sum(\gamma\lambda)^l\delta_{t+l}$. PPO заменяет сложный trust region из TRPO на clipped objective, который запрещает ratio $\pi_\theta/\pi_{old}$ уходить далеко от $1$ в выгодную для loss сторону.

\newpage
\section{DDPG и TD3: off-policy actor-critic для непрерывного управления}

\subsection*{Проблема непрерывных действий}
DQN требует максимум по действиям:
\[
    \max_a Q(s,a).
\]
В дискретном пространстве можно перебрать все действия. В непрерывном пространстве это дорого или невозможно. DDPG решает проблему, вводя детерминированного актора:
\[
    \mu_\theta(s)\approx \argmax_a Q_\phi(s,a).
\]

\subsection*{Deterministic Policy Gradient}
Для детерминированной политики $a=\mu_\theta(s)$ градиент:
\[
    \nabla_\theta J(\mu_\theta)=
    \frac1{1-\gamma}\E_{s\sim d^{\mu}}
    \left[\nabla_\theta\mu_\theta(s)\nabla_a Q^\mu(s,a)\big|_{a=\mu_\theta(s)}\right].
\]
Интуиция: critic говорит, в какую сторону в пространстве действий надо сдвинуть действие, а actor через chain rule меняет параметры, чтобы выдавать более выгодные действия.

\subsection*{DDPG}
DDPG - off-policy actor-critic с replay buffer и target networks. Есть:
\begin{itemize}
    \item actor $\mu_\theta(s)$;
    \item critic $Q_\phi(s,a)$;
    \item target actor $\mu_{\theta^-}$;
    \item target critic $Q_{\phi^-}$.
\end{itemize}

Critic target:
\[
    y=r+\gamma(1-done)Q_{\phi^-}(s',\mu_{\theta^-}(s')).
\]
Critic loss:
\[
    L_Q(\phi)=\E_{(s,a,r,s')\sim\cD}\left(Q_\phi(s,a)-y\right)^2.
\]
Actor objective:
\[
    J(\theta)=\E_{s\sim\cD}Q_\phi(s,
    \mu_\theta(s)).
\]
Actor update:
\[
    \nabla_\theta J\approx \E_{s\sim\cD}\left[
    \nabla_\theta\mu_\theta(s)\nabla_a Q_\phi(s,a)\big|_{a=\mu_\theta(s)}
    \right].
\]

Target networks обновляются мягко:
\[
    \phi^-\leftarrow \tau\phi+(1-\tau)\phi^-,
    \qquad
    \theta^-\leftarrow \tau\theta+(1-\tau)\theta^-.
\]

\subsection*{Exploration в DDPG}
Поскольку политика детерминированная, для исследования добавляют шум:
\[
    a_t=\mu_\theta(s_t)+\epsilon_t.
\]
Можно использовать гауссов шум
\[
    \epsilon_t\sim\mathcal N(0,\sigma^2I),
\]
или шум Орнштейна-Уленбека:
\[
    \epsilon_{t+1}=\alpha\epsilon_t+\mathcal N(0,\sigma^2I),
\]
который дает скоррелированные по времени возмущения. Это полезно в физических задачах, где случайное дерганье на каждом шаге неестественно.

\subsection*{Проблемы DDPG}
DDPG чувствителен к гиперпараметрам и ошибкам критика. Основные проблемы:
\begin{itemize}
    \item overestimation bias в $Q$;
    \item actor эксплуатирует ошибки критика;
    \item bootstrapping может разгонять ошибку;
    \item политика может стать нестабильной.
\end{itemize}
TD3 был предложен как набор исправлений DDPG.

\subsection*{TD3: Twin Delayed DDPG}
TD3 добавляет три ключевые идеи:
\begin{enumerate}
    \item twin critics;
    \item target policy smoothing;
    \item delayed policy updates.
\end{enumerate}

\subsection*{Twin critics и clipped double Q}
TD3 обучает два критика:
\[
    Q_{\phi_1}(s,a),\qquad Q_{\phi_2}(s,a).
\]
В таргете используется минимум:
\[
    y=r+\gamma(1-done)\min_{i=1,2}Q_{\phi_i^-}(s',a').
\]
Это уменьшает overestimation bias. Если один критик переоценил действие, минимум выбирает более осторожную оценку.

\subsection*{Target policy smoothing}
В таргетное действие добавляют сглаженный шум:
\[
    a'=\mu_{\theta^-}(s')+\epsilon,
\]
\[
    \epsilon\sim\clip(\mathcal N(0,\sigma^2),-c,c),
\]
\[
    a'\leftarrow\clip(a',a_{low},a_{high}).
\]
Таргет:
\[
    y=r+\gamma(1-done)\min_i Q_{\phi_i^-}(s',a').
\]
Интуиция: critic не должен давать слишком острые пики для отдельных действий. Сглаживание заставляет значение действия быть устойчивым к небольшим изменениям действия.

\subsection*{Delayed policy updates}
В TD3 actor обновляют реже, чем critic. Например, critic обновляется каждый gradient step, а actor и target networks - раз в $d$ шагов. Это нужно, потому что actor должен обучаться по более надежному критику. Если actor обновлять слишком часто, он начинает эксплуатировать временные ошибки $Q$.

\subsection*{Actor update в TD3}
Actor использует только один из критиков, например $Q_{\phi_1}$:
\[
    J(\theta)=\E_{s\sim\cD}Q_{\phi_1}(s,\mu_\theta(s)),
\]
\[
    \nabla_\theta J=\E_{s\sim\cD}\left[
    \nabla_\theta\mu_\theta(s)\nabla_aQ_{\phi_1}(s,a)|_{a=\mu_\theta(s)}
    \right].
\]

\subsection*{DDPG и TD3: сравнение}
\begin{center}
\begin{tabular}{p{4cm}p{5cm}p{5cm}}
\toprule
Компонент & DDPG & TD3 \\
\midrule
Critic & один $Q$ & два $Q$, минимум в таргете \\
Actor update & каждый шаг & с задержкой \\
Target action & $\mu^-(s')$ & $\mu^-(s')$ плюс clipped noise \\
Цель & непрерывное управление & более стабильный DDPG \\
\bottomrule
\end{tabular}
\end{center}

\subsection*{Что обязательно сказать на экзамене}
DDPG - это DQN-идея для непрерывных действий: вместо перебора $\argmax_a Q(s,a)$ вводится actor $\mu(s)$. Critic учится off-policy по replay buffer, actor обновляется через deterministic policy gradient. TD3 исправляет DDPG тремя трюками: два критика и минимум для борьбы с переоценкой, сглаживание таргетного действия и отложенные обновления actor.

\newpage
\section{Entropy RL и Soft Actor-Critic}

\subsection*{Maximum Entropy RL}
В обычном RL максимизируется ожидаемая сумма наград:
\[
    J(\pi)=\E_{\tau\sim\pi}\sum_{t\ge0}\gamma^t r_t.
\]
В Maximum Entropy RL добавляется энтропийный бонус:
\[
    J_{soft}(\pi)=\E_{\tau\sim\pi}\sum_{t\ge0}\gamma^t
    \left[r_t+\alpha H(\pi(\cdot\mid s_t))\right],
\]
где
\[
    H(\pi(\cdot\mid s))=-\E_{a\sim\pi}\log\pi(a\mid s).
\]
Эквивалентно:
\[
    J_{soft}(\pi)=\E_{\tau\sim\pi}\sum_{t\ge0}\gamma^t
    \left[r_t-\alpha\log\pi(a_t\mid s_t)\right].
\]
$\alpha$ - temperature. Чем больше $\alpha$, тем сильнее поощрение случайности.

\subsection*{Зачем нужна энтропия}
Энтропийный бонус:
\begin{itemize}
    \item улучшает exploration;
    \item мешает политике слишком рано схлопнуться в детерминированную;
    \item делает обучение устойчивее;
    \item позволяет получить стохастическую оптимальную политику.
\end{itemize}
Если награды в MDP всегда нулевые, оптимальная политика в max entropy смысле - политика максимальной энтропии, а не детерминированная.

\subsection*{Soft value functions}
Для простоты можно считать $\alpha=1$; в общем случае логарифм умножается на $\alpha$.

Soft Q-function:
\[
    Q_{soft}^\pi(s,a)=r(s,a)+\gamma\E_{s'}V_{soft}^\pi(s').
\]
Soft V-function:
\[
    V_{soft}^\pi(s)=\E_{a\sim\pi}\left[Q_{soft}^\pi(s,a)-\alpha\log\pi(a\mid s)\right].
\]
Soft Bellman equation:
\[
    Q_{soft}^\pi(s,a)=r(s,a)+\gamma\E_{s',a'}
    \left[Q_{soft}^\pi(s',a')-\alpha\log\pi(a'\mid s')\right].
\]

\subsection*{Soft policy improvement}
Для фиксированного soft critic политика улучшается через задачу:
\[
    \max_\pi \E_{a\sim\pi(\cdot\mid s)}\left[Q(s,a)-\alpha\log\pi(a\mid s)\right].
\]
Эквивалентно минимизации KL:
\[
    \pi_{new}=\argmin_\pi \KL\left(\pi(\cdot\mid s)\;\middle\|\;\frac{\exp(Q(s,\cdot)/\alpha)}{Z(s)}\right).
\]
Отсюда оптимальная soft-политика пропорциональна exponentiated Q:
\[
    \pi^*(a\mid s)\propto \exp(Q^*(s,a)/\alpha).
\]
В непрерывных действиях нормировочная константа обычно не считается явно, поэтому политику аппроксимируют нейросетью.

\subsection*{Soft Actor-Critic: общая идея}
SAC - off-policy actor-critic для max entropy RL. Он использует:
\begin{itemize}
    \item stochastic actor $\pi_\theta(a\mid s)$;
    \item два soft Q-критика $Q_{\phi_1},Q_{\phi_2}$;
    \item target Q-networks;
    \item replay buffer;
    \item entropy term в actor и critic targets.
\end{itemize}

Современная версия SAC обычно не использует отдельную $V_\psi$ сеть, хотя в ранней версии она была. В раннем варианте $V_\psi(s)$ аппроксимировала
\[
    V(s)=\E_{a\sim\pi}\left[Q(s,a)-\alpha\log\pi(a\mid s)\right]
\]
и использовалась для таргета критика.

\subsection*{Critic update в SAC}
Для перехода $(s,a,r,s',done)$ сэмплируем
\[
    a'\sim\pi_\theta(\cdot\mid s').
\]
Таргет:
\[
    y=r+\gamma(1-done)\left[
    \min_{i=1,2}Q_{\phi_i^-}(s',a')-
    \alpha\log\pi_\theta(a'\mid s')
    \right].
\]
Critic loss:
\[
    L_Q(\phi_i)=\E_{\cD}\left(Q_{\phi_i}(s,a)-y\right)^2,
    \qquad i=1,2.
\]
Минимум по двум критикам нужен, как в TD3, чтобы уменьшить overestimation bias.

\subsection*{Actor update}
Actor максимизирует soft value:
\[
    \max_\theta \E_{s\sim\cD,a\sim\pi_\theta}\left[
    \min_iQ_{\phi_i}(s,a)-\alpha\log\pi_\theta(a\mid s)
    \right].
\]
Эквивалентный loss для минимизации:
\[
    L_\pi(\theta)=\E_{s\sim\cD,\epsilon}\left[
    \alpha\log\pi_\theta(a_\theta(s,
    \epsilon)\mid s)-\min_iQ_{\phi_i}(s,a_\theta(s,\epsilon))
    \right].
\]
Здесь используется reparameterization trick:
\[
    a_\theta(s,\epsilon)=\tanh(\mu_\theta(s)+\sigma_\theta(s)\odot\epsilon),
    \qquad \epsilon\sim\mathcal N(0,I).
\]

\subsection*{Автоматическая настройка температуры}
Можно обучать $\alpha$, чтобы энтропия политики была близка к целевой $\bar H$:
\[
    L(\alpha)=\E_{a\sim\pi_\theta}\left[-\alpha(\log\pi_\theta(a\mid s)+\bar H)\right].
\]
Если энтропия слишком мала, $\alpha$ растет и сильнее поощряет случайность. Если энтропия слишком велика, $\alpha$ уменьшается.

\subsection*{SAC как развитие Soft Q-Learning}
Soft Q-Learning теоретически задает политику через
\[
    \pi(a\mid s)\propto\exp(Q(s,a)/\alpha).
\]
Но в непрерывных действиях постоянно сэмплировать из такого распределения трудно. SAC добавляет actor, который явно аппроксимирует эту soft-greedy политику. Это похоже на идею DDPG: вместо поиска argmax по $Q$ добавляем actor. Только в SAC actor стохастический и обучается с entropy regularization.

\subsection*{Что обязательно сказать на экзамене}
Maximum Entropy RL максимизирует не только награду, но и энтропию политики: $r-\alpha\log\pi(a\mid s)$. Soft Bellman equations содержат член $-\alpha\log\pi$. SAC - off-policy actor-critic: два Q-критика, стохастический actor, replay buffer, таргет $r+\gamma(\min Q-\alpha\log\pi)$. Actor оптимизирует $\alpha\log\pi-Q$, то есть балансирует высокую ценность и высокую энтропию.

\newpage
\section{Имитационное обучение, offline RL, Guided Cost Learning и GAIL}

\subsection*{Имитационное обучение}
В RL часто трудно задать функцию награды. Иногда проще получить демонстрации эксперта:
\[
    \cD_E=\{\tau_i\},\qquad \tau=(s_0,a_0,s_1,a_1,\dots).
\]
Цель imitation learning - научить агента вести себя как эксперт.

\subsection*{Behavioral Cloning}
Behavioral Cloning сводит задачу к supervised learning:
\[
    \max_\theta \sum_{(s,a)\in\cD_E}\log\pi_\theta(a\mid s).
\]
Для дискретных действий это cross-entropy classification, для непрерывных - регрессия или максимизация likelihood гауссовой политики.

Проблема: covariate shift. Модель обучалась на состояниях эксперта, но при собственной ошибке попадает в состояния, которых не было в датасете. Ошибки накапливаются.

\subsection*{DAgger}
DAgger борется с covariate shift:
\begin{enumerate}
    \item Обучить начальную политику на демонстрациях.
    \item Запустить текущую политику и собрать посещенные состояния.
    \item Попросить эксперта разметить правильные действия в этих состояниях.
    \item Добавить данные в датасет и переобучить политику.
\end{enumerate}
Минус: нужен эксперт, способный размечать новые состояния.

\subsection*{Offline RL}
Offline RL обучает политику по фиксированному датасету переходов:
\[
    \cD=\{(s,a,r,s')\},
\]
без новых взаимодействий со средой. Это важно, если реальные эксперименты дороги или опасны.

Главная проблема offline RL - distributional shift и extrapolation error. Если политика выбирает действия, которых почти нет в датасете, critic может сильно ошибаться:
\[
    Q(s,a)\quad \text{для out-of-distribution } a.
\]
Затем actor начинает эксплуатировать эти ошибки.

Типичные идеи offline RL:
\begin{itemize}
    \item ограничивать новую политику близостью к behavior policy;
    \item штрафовать OOD-действия;
    \item использовать conservative Q-learning;
    \item комбинировать RL objective с behavioral cloning loss.
\end{itemize}

\subsection*{Inverse Reinforcement Learning}
IRL пытается восстановить функцию награды, которую максимизирует эксперт. По демонстрациям нельзя однозначно восстановить награду: много разных reward functions могут объяснять одну и ту же политику. Поэтому нужны дополнительные предположения.

В Maximum Entropy IRL предполагают, что вероятность траектории эксперта пропорциональна экспоненте ее награды:
\[
    p_\theta(\tau)\propto \exp(R_\theta(\tau))\prod_t p(s_{t+1}\mid s_t,a_t),
\]
где
\[
    R_\theta(\tau)=\sum_t r_\theta(s_t,a_t).
\]
Нормировочная константа:
\[
    Z(\theta)=\int \exp(R_\theta(\tau))\prod_t p(s_{t+1}\mid s_t,a_t)\dd\tau.
\]

\subsection*{Guided Cost Learning}
Guided Cost Learning обучает cost/reward model по экспертным траекториям. Лог-правдоподобие экспертной траектории:
\[
    \log p_\theta(\tau)=R_\theta(\tau)-\log Z(\theta)+const.
\]
Градиент имеет вид разности признаков/наград на экспертных и сэмплированных траекториях:
\[
    \nabla_\theta \mathcal L
    =\E_{\tau\sim\cD_E}\nabla_\theta R_\theta(\tau)
    -\E_{\tau\sim p_\theta}\nabla_\theta R_\theta(\tau).
\]
Проблема: второе матожидание по текущей модели траекторий трудно считать. Поэтому используется сэмплирование траекторий текущей политики, которая обучается максимизировать найденную reward/cost function. Отсюда название guided: обучение reward и обучение policy направляют друг друга.

Общая схема:
\begin{enumerate}
    \item Имеем экспертные траектории.
    \item Обучаем reward/cost так, чтобы экспертные траектории были более вероятны, чем сгенерированные агентом.
    \item Обучаем политику на текущей reward/cost.
    \item Генерируем новые траектории агентом и повторяем.
\end{enumerate}

\subsection*{GAIL}
GAIL - Generative Adversarial Imitation Learning. Он избегает явного восстановления reward function и формулирует imitation как GAN-подобную игру.

Есть:
\begin{itemize}
    \item generator - политика $\pi_\theta$, порождающая траектории;
    \item discriminator $D_w(s,a)$, отличающий пары экспертных данных от пар агента.
\end{itemize}

Objective:
\[
    \min_\pi\max_D\;
    \E_{(s,a)\sim\pi_E}\log D(s,a)
    +\E_{(s,a)\sim\pi}\log(1-D(s,a))
    -\lambda H(\pi).
\]
Discriminator обучается как бинарный классификатор: эксперт - $1$, агент - $0$. Политика обучается так, чтобы обмануть дискриминатор.

\subsection*{Reward из дискриминатора}
Для обновления политики используют reward, полученный из discriminator. Возможные варианты:
\[
    r(s,a)=-\log(1-D(s,a))
\]
или
\[
    r(s,a)=\log D(s,a)-\log(1-D(s,a)).
\]
Затем политика обновляется обычным RL-методом, например TRPO/PPO.

\subsection*{Связь GAIL с occupancy measures}
Occupancy measure политики:
\[
    \rho_\pi(s,a)=\pi(a\mid s)\sum_{t=0}^{\infty}\gamma^t\Prob(s_t=s\mid\pi).
\]
GAIL можно понимать как минимизацию расстояния между occupancy measure эксперта и агента. Если агент посещает те же пары $(s,a)$ с теми же частотами, что эксперт, он имитирует поведение эксперта.

\subsection*{GAN и GAIL: отличие}
В GAN generator напрямую дифференцируемо преобразует шум в объект, и градиент от discriminator течет в generator. В GAIL generator - это политика внутри среды: действие влияет на будущие состояния через недифференцируемую динамику. Поэтому политика обучается не прямым backprop через discriminator, а RL-методом с reward от discriminator.

\subsection*{Что обязательно сказать на экзамене}
Behavioral cloning - supervised learning по демонстрациям, но страдает от covariate shift. Offline RL учится по фиксированному датасету и боится out-of-distribution действий. IRL восстанавливает награду из экспертных траекторий. Guided Cost Learning обучает reward через max entropy likelihood. GAIL заменяет явный reward adversarial-схемой: discriminator отличает эксперта от агента, а policy обучается обманывать discriminator через RL.

\newpage
\section{Многорукие бандиты: UCB и Thompson Sampling}

\subsection*{Постановка задачи}
В задаче многорукого бандита есть набор действий/ручек $a\in\cA$. На шаге $k$ агент выбирает $a_k$ и получает награду
\[
    r_k\sim p(r\mid a_k).
\]
Ценность действия:
\[
    Q(a)=\E[r\mid a].
\]
Оптимальная ценность:
\[
    V^*=\max_{a\in\cA}Q(a).
\]
Цель - минимизировать regret:
\[
    \mathcal R_T=TV^*-
    \sum_{k=0}^{T-1}Q(a_k).
\]
Эквивалентно:
\[
    \mathcal R_T=\sum_{a\in\cA}n_T(a)(V^*-Q(a)),
\]
где $n_T(a)$ - сколько раз выбрали действие $a$.

\subsection*{Exploration-exploitation dilemma}
Нужно балансировать:
\begin{itemize}
    \item exploitation - выбирать действие с лучшей текущей оценкой;
    \item exploration - пробовать плохо изученные действия, которые могут оказаться лучше.
\end{itemize}
Жадный алгоритм может навсегда застрять на неоптимальной ручке из-за случайно удачного начального сэмпла.

\subsection*{Оценка среднего}
Эмпирическая оценка:
\[
    Q_k(a)=\frac{\sum_{i=0}^{k}r_i\1\{a_i=a\}}{\sum_{i=0}^{k}\1\{a_i=a\}}.
\]
Инкрементальное обновление для выбранного действия:
\[
    Q_k(a_k)=Q_{k-1}(a_k)+\frac1{n_k(a_k)}\left(r_k-Q_{k-1}(a_k)\right).
\]
Для нестационарных бандитов используют постоянный шаг:
\[
    Q_k(a_k)=Q_{k-1}(a_k)+\alpha(r_k-Q_{k-1}(a_k)).
\]

\subsection*{$\varepsilon$-greedy}
Простой подход:
\[
    a_k=\begin{cases}
    \text{случайное действие},&\text{с вероятностью }\varepsilon,\\
    \argmax_a Q_k(a),&\text{с вероятностью }1-\varepsilon.
    \end{cases}
\]
При постоянном $\varepsilon$ regret обычно растет линейно, потому что агент продолжает с ненулевой вероятностью выбирать заведомо плохие действия.

\subsection*{Lai-Robbins lower bound}
Для хороших алгоритмов regret не может расти быстрее нижней логарифмической границы. В классической форме:
\[
    \E[\mathcal R_T]\ge
    \log T\sum_{a\ne a^*}
    \frac{V^*-Q(a)}{\KL(p(\cdot\mid a)\|p(\cdot\mid a^*))}.
\]
Алгоритм называется асимптотически оптимальным, если его regret растет логарифмически.

\subsection*{UCB: optimism in the face of uncertainty}
UCB строит верхнюю доверительную границу для каждого действия:
\[
    a_k=\argmax_a\left[Q_k(a)+U_k(a)\right].
\]
Идея: выбирать действие с лучшим оптимистичным значением. Если действие мало пробовали, $U_k(a)$ велик, и алгоритм его исследует. Если пробовали много, доверительный интервал сужается.

По неравенству Хёффдинга для наград в $[0,1]$:
\[
    \Prob(\mu\ge \hat\mu+u)\le e^{-2nu^2}.
\]
Отсюда доверительный бонус:
\[
    U_k(a)=\sqrt{\frac{-\log\delta}{2n_k(a)}}.
\]
В практическом UCB часто используют:
\[
    a_k=\argmax_a\left[Q_k(a)+c\sqrt{\frac{\log k}{n_k(a)}}\right].
\]

\subsection*{Алгоритм UCB}
\begin{enumerate}
    \item Каждое действие выбрать хотя бы один раз.
    \item На шаге $k$ выбрать
    \[
        a_k=\argmax_a\left[Q_k(a)+c\sqrt{\frac{\log k}{n_k(a)}}\right].
    \]
    \item Получить $r_k$.
    \item Обновить $Q_k(a_k)$ и $n_k(a_k)$.
\end{enumerate}

UCB не требует случайного exploration-параметра $\varepsilon$; исследование возникает из неопределенности.

\subsection*{Thompson Sampling}
Thompson Sampling - байесовский probability matching. Для каждого действия есть неизвестный параметр $\theta_a$ распределения наград. Поддерживаем posterior:
\[
    p(\theta_a\mid \text{data}).
\]
На каждом шаге:
\begin{enumerate}
    \item Для каждого действия сэмплировать
    \[
        \tilde\theta_a\sim p(\theta_a\mid\text{data}).
    \]
    \item Посчитать сэмплированную ценность
    \[
        \tilde Q(a)=\E[r\mid \tilde\theta_a].
    \]
    \item Выбрать
    \[
        a_k=\argmax_a \tilde Q(a).
    \]
    \item Получить награду и обновить posterior выбранного действия.
\end{enumerate}

Действие выбирается с вероятностью, равной posterior-вероятности быть оптимальным:
\[
    \pi(a)=\Prob\left(Q(a)=\max_{\hat a}Q(\hat a)\mid\text{data}\right).
\]

\subsection*{Bernoulli bandit и Beta posterior}
Если награда $r\in\{0,1\}$:
\[
    r\sim \mathrm{Bernoulli}(\theta_a).
\]
Сопряженный prior:
\[
    \theta_a\sim\mathrm{Beta}(\alpha_a,\beta_a).
\]
После наблюдения $r$:
\[
    \alpha_a\leftarrow \alpha_a+r,
    \qquad
    \beta_a\leftarrow \beta_a+1-r.
\]
Среднее posterior:
\[
    \E[\theta_a]=\frac{\alpha_a}{\alpha_a+
    \beta_a}.
\]
Но Thompson Sampling выбирает не по среднему posterior, а сэмплирует $\tilde\theta_a\sim\mathrm{Beta}(\alpha_a,\beta_a)$ и выбирает максимум.

\subsection*{UCB против Thompson Sampling}
UCB использует оптимизм через верхнюю границу доверительного интервала. Thompson Sampling использует posterior uncertainty: если действие плохо изучено, его posterior широк, и оно иногда сэмплируется как лучшее. Оба метода естественно уменьшают exploration по мере накопления данных.

\subsection*{Что обязательно сказать на экзамене}
Бандит - MDP без состояния, где нужно минимизировать regret. UCB выбирает действие по $Q_k(a)+c\sqrt{\log k/n_k(a)}$, реализуя оптимизм перед неопределенностью. Thompson Sampling поддерживает posterior на параметры наград и выбирает действие, которое оказалось лучшим в сэмпле из posterior. Для Bernoulli-бандита Beta prior обновляется как $\alpha\leftarrow\alpha+r$, $\beta\leftarrow\beta+1-r$.

\newpage
\section{Monte Carlo Tree Search}

\subsection*{Планирование}
Планирование отличается от обучения политики тем, что для текущего состояния $s_0$ мы ищем хороший план действий:
\[
    a_0,a_1,a_2,\dots
\]
Если известна модель переходов и наград, задача:
\[
    \argmax_{(a_0,a_1,\dots)}
    \E\left[R(\tau)\mid s_0,a_0,a_1,\dots\right].
\]
В дискретных действиях можно строить дерево возможных продолжений. Но полное дерево растет экспоненциально по глубине:
\[
    |\cA|^H.
\]
MCTS строит дерево выборочно, концентрируясь на перспективных ветвях.

\subsection*{Дерево MDP}
Корень дерева соответствует текущему состоянию $s_0$. Ребро соответствует действию $a$. Узел на глубине $t$ соответствует плану $(a_0,\dots,a_{t-1})$. В узлах/ребрах хранят статистику:
\begin{itemize}
    \item $n(N,a)$ - сколько раз проходили по ребру $(N,a)$;
    \item $Q(N,a)$ - оценка ценности действия $a$ из узла $N$;
    \item $n(N)=\sum_a n(N,a)$ - число посещений узла.
\end{itemize}

\subsection*{Четыре фазы MCTS}
Одна итерация MCTS состоит из:
\begin{enumerate}
    \item Selection - спуск по уже построенному дереву.
    \item Expansion - добавление новых узлов.
    \item Simulation/Evaluation - оценка новых листьев.
    \item Backpropagation/Update - обновление статистик по пройденному пути.
\end{enumerate}

\subsection*{Selection}
Selection выбирает путь от корня до листа по tree policy. Обычно используют UCB-подобное правило:
\[
    a=\argmax_{a\in\cA}\left[
    Q(N,a)+C\sqrt{\frac{\log n(N)}{n(N,a)}}
    \right].
\]
Первое слагаемое - exploitation: идти в ветви с высокой оценкой. Второе - exploration: пробовать редко посещенные ветви.

При спуске через ребро $(N,a)$ симулятор генерирует следующий переход:
\[
    s',r\sim p(s',r\mid s,a).
\]

\subsection*{Expansion}
Когда дошли до листа, дерево расширяют: для возможных действий создают новые дочерние узлы. Если действий слишком много, можно расширять не все действия, а случайный поднабор или progressive widening.

\subsection*{Simulation / Evaluation}
Нужно оценить ценность нового листа. В классическом MCTS запускают rollout до конца эпизода с default policy, часто случайной:
\[
    \pi_{default}(a\mid s).
\]
Получают Monte-Carlo оценку возврата:
\[
    \hat R=\sum_{t=0}^{T}\gamma^t r_t.
\]
В современных вариантах вместо случайного rollout можно использовать value network:
\[
    \hat V(s_{leaf})\approx V(s_{leaf}).
\]

\subsection*{Backpropagation}
После simulation полученную оценку возвращают назад по пройденному пути. Для каждого ребра $(N,a)$:
\[
    n(N,a)\leftarrow n(N,a)+1,
\]
\[
    Q(N,a)\leftarrow Q(N,a)+\frac1{n(N,a)}\left(\hat R-Q(N,a)\right).
\]
Если в пути были промежуточные награды, для каждого ребра используют соответствующий reward-to-go после этого ребра.

\subsection*{Итоговая политика из MCTS}
После $K$ итераций MCTS выбирает действие в корне. Можно взять:
\[
    a_0=\argmax_a n(N_0,a)
\]
или задать стохастическую политику по посещениям:
\[
    \pi(a\mid s_0)\propto n(N_0,a)^{1/T},
\]
где $T$ - температура. При $T\to0$ выбор становится жадным.

\subsection*{Использование в реальной среде}
Обычно MCTS работает online:
\begin{enumerate}
    \item В текущем состоянии $s_t$ сделать много итераций MCTS.
    \item Выбрать реальное действие $a_t$.
    \item Среда переходит в $s_{t+1}$.
    \item Сохранить поддерево, соответствующее выбранному действию и новому состоянию.
    \item Повторить планирование из нового корня.
\end{enumerate}
Это похоже на receding horizon planning: постоянно перепланируем.

\subsection*{MCTS и бандиты}
В каждом узле MCTS нужно выбрать, какую ветвь исследовать. Это локальная задача многорукого бандита. Поэтому UCB естественно используется как tree policy. MCTS можно понимать как много вложенных бандитных задач, по одной в каждом узле дерева.

\subsection*{Плюсы и минусы MCTS}
Плюсы:
\begin{itemize}
    \item не строит полное дерево;
    \item концентрируется на перспективных ветвях;
    \item anytime algorithm: чем больше итераций, тем лучше решение;
    \item хорошо работает в дискретных action spaces и играх.
\end{itemize}
Минусы:
\begin{itemize}
    \item требует симулятор/модель;
    \item дорог по времени на каждый реальный action;
    \item сложен для непрерывных действий без модификаций;
    \item качество зависит от rollout policy/value heuristic.
\end{itemize}

\subsection*{AlphaZero-style MCTS}
В современных методах MCTS комбинируют с нейросетями:
\begin{itemize}
    \item policy network дает prior $P(s,a)$ для действий;
    \item value network оценивает лист вместо случайного rollout;
    \item tree policy использует PUCT:
    \[
        a=\argmax_a\left[Q(s,a)+cP(s,a)\frac{\sqrt{N(s)}}{1+N(s,a)}\right].
    \]
\end{itemize}
Это делает поиск намного эффективнее.

\subsection*{Что обязательно сказать на экзамене}
MCTS - метод планирования с симулятором для дискретных действий. Одна итерация: selection по UCB, expansion, simulation/evaluation, backpropagation. В ребрах хранят $Q(N,a)$ и $n(N,a)$. После многих итераций действие выбирают по максимальному числу посещений или по распределению посещений.

\newpage
\section{LQR, iLQR и общая схема Model-based RL}

\subsection*{Model-based RL}
Model-based RL явно или неявно обучает модель динамики среды:
\[
    p_\theta(s'\mid s,a)
\]
или детерминированную модель
\[
    s_{t+1}=f_\theta(s_t,a_t).
\]
Общая схема:
\begin{enumerate}
    \item Собрать переходы $(s,a,r,s')$ в датасет.
    \item Обучить модель динамики и, если нужно, модель награды.
    \item Использовать модель для планирования или генерации ``сновидений''.
    \item Улучшить политику.
    \item Собрать новые данные и повторить.
\end{enumerate}

Плюсы model-based RL:
\begin{itemize}
    \item высокая sample efficiency;
    \item можно планировать без реального риска;
    \item можно использовать знания физики.
\end{itemize}
Минусы:
\begin{itemize}
    \item модель ошибается;
    \item планировщик может эксплуатировать ошибки модели;
    \item ошибки накапливаются при длинных rollout.
\end{itemize}

\subsection*{Планирование в непрерывном управлении}
Если динамика и награда дифференцируемы:
\[
    s_{t+1}=f(s_t,a_t),
    \qquad
    r_t=r(s_t,a_t),
\]
можно оптимизировать последовательность действий:
\[
    \max_{a_{0:T-1}}\sum_{t=0}^{T-1}r(s_t,a_t)
\]
при ограничениях динамики. Прямой способ - backpropagation through time по действиям. Но для систем управления есть классические методы: LQR и iLQR.

\subsection*{Linear Quadratic Regulator}
LQR решает задачу с линейной динамикой и квадратичной стоимостью. Обычно формулируется как минимизация cost:
\[
    x_{t+1}=Ax_t+Bu_t,
\]
\[
    J=\sum_{t=0}^{T-1}(x_t^TQx_t+u_t^TRu_t)+x_T^TQ_f x_T.
\]
Здесь $Q\succeq0$, $R\succ0$. Нужно найти управляющие воздействия $u_t$, минимизирующие $J$.

\subsection*{Value function в LQR}
Для LQR оптимальная value function квадратична:
\[
    V_t(x)=x^TP_t x.
\]
На последнем шаге:
\[
    P_T=Q_f.
\]
Рекурсия Риккати назад во времени:
\[
    P_t=Q+A^TP_{t+1}A
    -A^TP_{t+1}B(R+B^TP_{t+1}B)^{-1}B^TP_{t+1}A.
\]
Оптимальное управление линейно по состоянию:
\[
    u_t=-K_t x_t,
\]
где
\[
    K_t=(R+B^TP_{t+1}B)^{-1}B^TP_{t+1}A.
\]

\subsection*{Бесконечный горизонт}
Для стационарной задачи бесконечного горизонта $P_t\to P$, где $P$ решает алгебраическое уравнение Риккати:
\[
    P=Q+A^TPA-A^TPB(R+B^TPB)^{-1}B^TPA.
\]
Тогда управление стационарно:
\[
    u=-Kx.
\]

\subsection*{Интуиция LQR}
LQR балансирует две цели:
\begin{itemize}
    \item держать состояние близко к нулю через штраф $x^TQx$;
    \item не тратить слишком большое управление через штраф $u^TRu$.
\end{itemize}
Если $R$ большой, управление осторожное. Если $Q$ большой, система агрессивнее возвращает состояние к цели.

\subsection*{iLQR: iterative LQR}
Реальные системы часто нелинейны:
\[
    x_{t+1}=f(x_t,u_t),
\]
а cost не обязательно квадратичный:
\[
    J=\sum_t \ell(x_t,u_t)+\ell_f(x_T).
\]
iLQR решает задачу итеративно: вокруг текущей номинальной траектории линеаризует динамику и квадратично аппроксимирует cost, затем решает локальную LQR-задачу.

\subsection*{Линеаризация динамики}
Пусть есть номинальная траектория $(\bar x_t,\bar u_t)$. Введем отклонения:
\[
    \delta x_t=x_t-\bar x_t,
    \qquad
    \delta u_t=u_t-\bar u_t.
\]
Динамика приближенно:
\[
    \delta x_{t+1}\approx f_x(t)\delta x_t+f_u(t)\delta u_t,
\]
где
\[
    f_x(t)=\frac{\partial f}{\partial x}\bigg|_{\bar x_t,\bar u_t},
    \qquad
    f_u(t)=\frac{\partial f}{\partial u}\bigg|_{\bar x_t,\bar u_t}.
\]

\subsection*{Квадратичная аппроксимация cost}
Cost раскладывают до второго порядка:
\[
    \ell(x_t,u_t)\approx \ell_t+
    \ell_x^T\delta x_t+\ell_u^T\delta u_t+
    \frac12\delta x_t^T\ell_{xx}\delta x_t+
    \delta u_t^T\ell_{ux}\delta x_t+
    \frac12\delta u_t^T\ell_{uu}\delta u_t.
\]
После этого локальная задача становится LQR-подобной.

\subsection*{Backward pass iLQR}
Определяют локальную Q-функцию для отклонений:
\[
    Q(\delta x,
    \delta u)=\ell(x,u)+V_{t+1}(f(x,u)).
\]
Ее производные дают оптимальное локальное управление:
\[
    \delta u_t^*=k_t+K_t\delta x_t,
\]
где
\[
    k_t=-Q_{uu}^{-1}Q_u,
    \qquad
    K_t=-Q_{uu}^{-1}Q_{ux}.
\]
Здесь $k_t$ - feedforward correction, $K_t$ - feedback gain.

\subsection*{Forward pass iLQR}
После backward pass выполняют rollout новой траектории:
\[
    u_t^{new}=\bar u_t+\alpha k_t+K_t(x_t^{new}-\bar x_t),
\]
где $\alpha\in(0,1]$ выбирается line search. Если cost уменьшился, новая траектория принимается. Затем процесс повторяется.

\subsection*{Алгоритм iLQR}
\begin{enumerate}
    \item Инициализировать последовательность управлений $u_{0:T-1}$.
    \item Сделать forward rollout через настоящую динамику и получить номинальную траекторию.
    \item Линеаризовать динамику и квадратично аппроксимировать cost вдоль траектории.
    \item Backward pass: посчитать $k_t,K_t$.
    \item Forward pass: применить новое управление с line search.
    \item Повторять до сходимости.
\end{enumerate}

\subsection*{Связь iLQR с model-based RL}
Если настоящая динамика неизвестна, можно обучить модель $f_\theta(s,a)$ и применять iLQR внутри этой модели. Тогда получаем model-based RL pipeline:
\begin{enumerate}
    \item собрали данные на реальной среде;
    \item обучили дифференцируемую модель динамики;
    \item планируем через iLQR/CEM/MPC;
    \item исполняем первое действие;
    \item добавляем новые данные и переобучаем модель.
\end{enumerate}

\subsection*{Шумная динамика}
Если переход стохастический:
\[
    s_{t+1}=f(s_t,a_t)+\epsilon_t,
\]
то планирование по среднему может быть недостаточно. Возможные подходы:
\begin{itemize}
    \item оптимизировать ожидаемый cost;
    \item учитывать ковариацию шума;
    \item использовать ансамбли моделей для неопределенности;
    \item планировать робастно или риск-чувствительно;
    \item перепланировать на каждом шаге в MPC-стиле.
\end{itemize}

\subsection*{Что обязательно сказать на экзамене}
LQR решает линейно-квадратичную задачу: линейная динамика, квадратичный cost, оптимальное управление $u=-Kx$, матрицы находятся рекурсией Риккати. iLQR обобщает это на нелинейную динамику: вокруг текущей траектории линеаризует динамику, квадратично аппроксимирует cost, решает локальный LQR, делает forward rollout и повторяет. В model-based RL модель динамики можно обучить по данным и использовать LQR/iLQR/MPC для планирования.

\end{document}
